% ============================================================
% 第五章  整体类域论: 主要定理的陈述
% 源: Milne CFT v4.03, Chapter V (书页 p.147--177)
% OCR: 行 5537--6795; 脚注 1--13 从原书 PDF 恢复
% 勘误: p.154-155 三处 1-s -> s-1 (Yiming Tang); 定理 3.6 补 "L/K 有模 m" (Zhuo Zhang)
% 标签前缀: v.  公式 (43)(44) 用 \tag 硬编码
% ============================================================
\chapter{整体类域论: 主要定理的陈述}\label{ch:V}
\index{zhengtileiyulun@整体类域论(global class field theory)}

\begin{quote}
\itshape La théorie du corps de classes a une réputation de difficulté qui est en partie justifiée. Mais il faut faire une distinction: il n'est peut-être pas en effet dans la science de théorie où tout à la fois les démonstrations soient aussi ardues, et les résultats d'une aussi parfaite simplicité et d'une aussi grande puissance.
\end{quote}
\hspace*{\fill}J.~Herbrand, 1936, p.~2\footnote{类域论以困难著称, 这在一定程度上是有道理的. 但必须作出区分: 科学中也许再没有别的理论, 其证明如此艰深, 而其结果却又如此完美地简洁、如此威力巨大. (原书脚注给出的英译, 此处译出.)}

本章陈述并解释整体类域论的主要定理. 这些定理将在第 \ref{ch:VII} 章中证明.

本章始终设 $K$ 为数域, 尽管大多数结果对 $\F_p(T)$ 的有限扩张同样成立.

回顾: 对数域 $K$, 我们把 $K$ 的\textbf{素} (prime) 定义为 $K$ 的非平凡绝对值的等价类. 素有两类: \textbf{有限素} (finite prime), 可与 $\mathcal{O}_K$ 的素理想等同; 以及\textbf{无限素} (infinite prime). 实无限素可与一个嵌入 $K\hookrightarrow\R$ 等同, 复无限素可与一对共轭嵌入 $K\hookrightarrow\CC$ 等同. 我们用 $\mathfrak{p}$ 或 $v$ 表示素 (有限或无限), 用 $S$ 表示 $K$ 的一个有限素集, 同时也表示 $K$ 的有限扩张 $L$ 中位于其上的素集. 无限素集记作 $S_\infty$.

$K$ 在素 $\mathfrak{p}$ (相应地 $v$) 处的完备化记作 $K_{\mathfrak{p}}$ (相应地 $K_v$), 包含映射 $K\hookrightarrow K_{\mathfrak{p}}$ (相应地 $K\hookrightarrow K_v$) 记作 $a\mapsto a_{\mathfrak{p}}$ (相应地 $a\mapsto a_v$).

\section{射线类群}\label{sec:V.1}

\subsection*{与 $S$ 互素的理想}

设 $I=I_K$ 为 $K$ 的分式理想群. 对 $K$ 的素的有限集 $S$, 定义 $I^S$ 为由不在 $S$ 中的素理想生成的 $I$ 的子群. 每个 $\mathfrak{a}\in I^S$ 有唯一的分解
\[
\mathfrak{a}=\mathfrak{p}_1^{n_1}\cdots\mathfrak{p}_s^{n_s},\quad\mathfrak{p}_i\notin S,\quad n_i\in\Z,
\]
因此 $I^S$ 可与由不在 $S$ 中的素理想生成的自由 Abel 群等同. 定义
\[
K^S=\{a\in K^{\times}\mid(a)\in I^S\}=\{a\in K^{\times}\mid\mathrm{ord}_{\mathfrak{p}}(a)=0\ \text{对所有有限}\ \mathfrak{p}\in S\}.
\]
设 $i\colon K^S\to I^S$ 为把 $a\in K^S$ 送到理想 $a\mathcal{O}_K$ 的映射.

例如, 当 $K=\Q$ 且 $S$ 为整除整数 $n$ 的素数集时, $I^S$ 是分式理想
\[
\{(r/s)\mid r,s\in\Z,\ \gcd(r,n)=1=\gcd(s,n)\}
\]
的集合, 而
\[
\Q^S=\{r/s\mid r,s\in\Z,\ \gcd(r,n)=1=\gcd(s,n)\}.
\]
此时自然映射 $\Q^S\to I^S$ 是满射, 核为 $\{\pm1\}$.

\begin{lemma}\label{lem:v.1.1}
对 $\mathcal{O}_K$ 中素理想的每个有限集 $S$, 序列
\[
0\to U_K\to K^S\to I^S\to C\to0
\]
正合. (此处 $U_K=\mathcal{O}_K^{\times}$, 而 $C$ 是完整理想类群 $I/i(K^{\times})$.)
\end{lemma}

\begin{proof}
为证 $I^S\to C$ 满, 须证每个理想类 $\mathcal{C}$ 都由 $I^S$ 中的理想代表. 设 $\mathfrak{a}$ 代表 $\mathcal{C}$. 则 $\mathfrak{a}=\mathfrak{b}\mathfrak{c}^{-1}$, 其中 $\mathfrak{b},\mathfrak{c}$ 为整理想, 且对任意 $c\in\mathfrak{c}$, $\mathfrak{a}\cdot(c)=\mathfrak{b}\mathfrak{c}^{-1}\cdot(c)$ 是整理想, 故可设 $\mathfrak{a}$ 本身是整理想. 写 $\mathfrak{a}=\prod_{\mathfrak{p}\in S}\mathfrak{p}^{n(\mathfrak{p})}\mathfrak{b}$, 其中 $\mathfrak{b}\in I^S$. 对每个 $\mathfrak{p}\in S$, 取 $\pi_{\mathfrak{p}}\in\mathfrak{p}\setminus\mathfrak{p}^2$, 使 $\mathrm{ord}_{\mathfrak{p}}(\pi_{\mathfrak{p}})=1$. 由中国剩余定理, 存在 $a\in\mathcal{O}_K$ 使得
\[
a\equiv\pi_{\mathfrak{p}}^{n(\mathfrak{p})}\bmod\mathfrak{p}^{n(\mathfrak{p})+1}\qquad(\text{对所有 }\mathfrak{p}\in S).
\]
这些同余式蕴含 $\mathrm{ord}_{\mathfrak{p}}(a)=n(\mathfrak{p})$ 对所有 $\mathfrak{p}\in S$ 成立, 于是 $(a)=\prod_{\mathfrak{p}\in S}\mathfrak{p}^{n(\mathfrak{p})}\mathfrak{b}'$, 其中 $\mathfrak{b}'\in I^S$. 现在 $a^{-1}\mathfrak{a}\in I^S$ 且与 $\mathfrak{a}$ 在 $C$ 中代表同一个类.

其次, 若 $\mathfrak{a}\in I^S$ 在 $C$ 中映为零, 则 $\mathfrak{a}=(\alpha)$ 对某个 $\alpha\in K^S$ 成立, 且 $\alpha$ 在相差一个单位的意义下唯一确定. 这就证明了其余各处的正合性.
\end{proof}

\begin{remark}\label{rem:v.1.2}
事实上, $C$ 中每个类都可由 $I^S$ 中的一个整理想 $\mathfrak{a}$ 代表: 设该类由 $\mathfrak{a}\in I^S$ 代表; 写 $\mathfrak{a}=\mathfrak{b}\mathfrak{c}^{-1}$, 其中 $\mathfrak{b},\mathfrak{c}$ 为 $I^S$ 中的整理想, 取非零的 $c\in\mathfrak{c}\cap K^S$ (由中国剩余定理存在), 并注意 $c\mathfrak{a}$ 是整理想.
\end{remark}

\subsection*{模}

\begin{definition}\label{def:v.1.3}
$K$ 的一个\emph{模} (modulus) 是一个函数
\[
m\colon\{K\ \text{的素}\}\to\Z
\]
满足:
\begin{enumerate}[(a)]
\item $m(\mathfrak{p})\ge 0$ 对所有素 $\mathfrak{p}$, 且除有限多个 $\mathfrak{p}$ 外 $m(\mathfrak{p})=0$;
\item 若 $\mathfrak{p}$ 为实素, 则 $m(\mathfrak{p})=0$ 或 $1$;
\item 若 $\mathfrak{p}$ 为复素, 则 $m(\mathfrak{p})=0$.
\end{enumerate}
\end{definition}

习惯上记
\[
\mathfrak{m}=\prod_{\mathfrak{p}}\mathfrak{p}^{m(\mathfrak{p})}.
\]
称模 $\mathfrak{m}=\prod\mathfrak{p}^{m(\mathfrak{p})}$ \textbf{整除}模 $\mathfrak{n}=\prod_{\mathfrak{p}}\mathfrak{p}^{n(\mathfrak{p})}$, 若 $m(\mathfrak{p})\le n(\mathfrak{p})$ 对所有 $\mathfrak{p}$ 成立. 特别地, 素 $\mathfrak{p}$ 整除模 $\mathfrak{m}$ 当且仅当 $m(\mathfrak{p})>0$.

模 $\mathfrak{m}$ 可写成
\[
\mathfrak{m}=\mathfrak{m}_\infty\mathfrak{m}_0,
\]
其中 $\mathfrak{m}_\infty$ 是一些实素之积, $\mathfrak{m}_0$ 是素理想的正数次幂之积, 故可与 $\mathcal{O}_K$ 中的一个理想等同.

\subsection*{射线类群}

对模 $\mathfrak{m}$, 定义 $K_{\mathfrak{m},1}$ 为满足如下条件的 $a\in K^{\times}$ 的集合:
\[
\begin{cases}
\mathrm{ord}_{\mathfrak{p}}(a-1)\ge m(\mathfrak{p}) & \text{对所有整除 }\mathfrak{m}\text{ 的有限 }\mathfrak{p},\\[2pt]
a_{\mathfrak{p}}>0 & \text{对所有整除 }\mathfrak{m}\text{ 的实 }\mathfrak{p}.
\end{cases}
\]
注意
\[
\mathrm{ord}_{\mathfrak{p}}(a-1)\ge m(\mathfrak{p})
\iff \pi^{m(\mathfrak{p})}\mid(a_{\mathfrak{p}}-1)
\iff a\mapsto 1\ \text{于}\ (\mathcal{O}_{\mathfrak{p}}/\mathfrak{p}^{m(\mathfrak{p})})^{\times}\simeq(\hat{\mathcal{O}}_{\mathfrak{p}}/\hat{\mathfrak{p}}^{m(\mathfrak{p})})^{\times},
\]
其中 $\pi$ 是 $K$ 在 $\mathfrak{p}$ 处的完备化 $K_{\mathfrak{p}}$ 中的素元. 令
\[
S(\mathfrak{m})=\{\text{整除 }\mathfrak{m}\text{ 的素}\}.
\]
对任意 $a\in K_{\mathfrak{m},1}$ 及整除 $\mathfrak{m}$ 的素理想 $\mathfrak{p}$, 有 $\mathrm{ord}_{\mathfrak{p}}(a-1)>0=\mathrm{ord}_{\mathfrak{p}}(1)$, 故
\[
\mathrm{ord}_{\mathfrak{p}}(a)=\mathrm{ord}_{\mathfrak{p}}\bigl((a-1)+1\bigr)=0.
\]
因此对任意 $a\in K_{\mathfrak{m},1}$, 理想 $(a)$ 属于 $I^{S(\mathfrak{m})}$. 设 $i$ 表示映射 $a\mapsto(a)\colon K_{\mathfrak{m},1}\to I^{S(\mathfrak{m})}$. 商群
\[
C_{\mathfrak{m}}=I^{S(\mathfrak{m})}/i(K_{\mathfrak{m},1})
\]
称为\textbf{模 $\mathfrak{m}$ 的 (射线) 类群}\index{shexianleiqun@射线类群(ray class group)}\index{leiqun@类群(class group)!shexianleiqun@射线类群(ray class group)}.

\begin{example}\label{ex:v.1.4}
表达式 $\mathfrak{m}=(2)^3\cdot(17)^2\cdot(19)\cdot\infty$ 是 $\Q$ 的一个模, 其中 $\mathfrak{m}_0=(2)^3\cdot(17)^2\cdot(19)$, $\mathfrak{m}_\infty=\infty$ (此处 $\infty$ 表示 $\Q$ 的唯一无限素). 此外, $\Q_{\mathfrak{m},1}$ 由满足
\[
\begin{cases}
\mathrm{ord}_2(a-1)\ge 3,\\
\mathrm{ord}_{17}(a-1)\ge 2,\\
\mathrm{ord}_{19}(a-1)\ge 1
\end{cases}
\]
的正有理数 $a$ 组成. 条件 ``$2$ 处'' 是说 $a$ 是两个奇整数之商 $a=b/c$, 且 $bc^{-1}$ 在 $(\Z/8\Z)^{\times}$ 中的像为 $1$. 其余条件可类似表述.
\end{example}

\begin{lemma}\label{lem:v.1.5}
设 $S$ 为 $K$ 的素理想的有限集. 则每个 $\alpha\in K^S$ 可写成 $\alpha=a/b$, 其中 $a,b\in\mathcal{O}_K\cap K^S$.
\end{lemma}

\begin{proof}
因 $\alpha\in K^S$, 有 $(\alpha)=\mathfrak{a}/\mathfrak{b}$, 其中 $\mathfrak{a},\mathfrak{b}$ 为 $I^S$ 中的整理想. 显然 $\mathfrak{a}$ 与 $\mathfrak{b}$ 代表理想类群中同一个元 $C$, 由注记 \ref{rem:v.1.2} 可取 $I^S$ 中的整理想 $\mathfrak{c}$ 代表 $C^{-1}$. 于是 $(\alpha)=\mathfrak{ac}/\mathfrak{bc}=(a)/(b)$, 对某个 $a,b\in\mathcal{O}_K\cap K^S$ 成立.
\end{proof}

\begin{proposition}\label{prop:v.1.6}
$C_{\mathfrak{m}}$ 中每个类都由一个整理想 $\mathfrak{a}$ 代表, 且两个整理想 $\mathfrak{a}$ 与 $\mathfrak{b}$ 代表 $C_{\mathfrak{m}}$ 中同一个类, 当且仅当存在非零的 $a,b\in\mathcal{O}_K$ 使得 $a\mathfrak{a}=b\mathfrak{b}$,
\[
a\equiv b\equiv 1\bmod\mathfrak{m}_0
\]
且 $a$ 与 $b$ 在每个整除 $\mathfrak{m}$ 的实素处有相同的符号.
\end{proposition}

\begin{proof}
设该类由 $\mathfrak{a}\in I^S$ 代表, 写 $\mathfrak{a}=\mathfrak{b}\mathfrak{c}^{-1}$, 其中 $\mathfrak{b},\mathfrak{c}$ 为 $I^S$ 中的整理想. 中国剩余定理表明存在非零的 $c\in\mathfrak{c}\cap K_{\mathfrak{m}_0,1}$, 而强逼近定理表明可取 $c$ 在所有实素处均为正. 此时 $c\mathfrak{a}$ 是整理想且与 $\mathfrak{a}$ 在 $C_{\mathfrak{m}}$ 中代表同一个类. 结论的第二部分由引理 \ref{lem:v.1.5} 推出.
\end{proof}

于是, 例如, 通常的理想类群可与下述等价关系下整理想的集合等同: $\mathfrak{a}\sim\mathfrak{b}$ 当且仅当 $a\mathfrak{a}=b\mathfrak{b}$ 对某对非零的 $a,b\in\mathcal{O}_K$ 成立.

\begin{theorem}\label{thm:v.1.7}
对 $K$ 的每个模 $\mathfrak{m}$, 有正合列
\[
0\to U/U_{\mathfrak{m},1}\to K_{\mathfrak{m}}/K_{\mathfrak{m},1}\to C_{\mathfrak{m}}\to C\to0
\]
与典范同构
\[
K_{\mathfrak{m}}/K_{\mathfrak{m},1}\simeq\prod_{\substack{\mathfrak{p}\ \mathrm{real}\\ \mathfrak{p}|\mathfrak{m}}}\{\pm\}\times\prod_{\substack{\mathfrak{p}\ \mathrm{finite}\\ \mathfrak{p}|\mathfrak{m}}}(\mathcal{O}_K/\mathfrak{p}^{m(\mathfrak{p})})^{\times}
\simeq\prod_{\substack{\mathfrak{p}\ \mathrm{real}\\ \mathfrak{p}|\mathfrak{m}}}\{\pm\}\times(\mathcal{O}_K/\mathfrak{m}_0)^{\times},
\]
其中
\begin{align*}
K_{\mathfrak{m}}&=K^{S(\mathfrak{m})}=\{\alpha\in K^{\times}\mid \mathrm{ord}_{\mathfrak{p}}(\alpha)=0\ \text{对所有 }\mathfrak{p}|\mathfrak{m}_0\},\\
U&=\mathcal{O}_K^{\times}\quad(\text{$K$ 的单位群}),\\
U_{\mathfrak{m},1}&=U\cap K_{\mathfrak{m},1}.
\end{align*}
因此 $C_{\mathfrak{m}}$ 是有限群, 其阶为
\[
h_{\mathfrak{m}}=h\cdot(U:U_{\mathfrak{m},1})^{-1}\cdot 2^{r_0}\cdot\Nm(\mathfrak{m}_0)\cdot\prod_{\mathfrak{p}|\mathfrak{m}_0}\Bigl(1-\frac{1}{\Nm\mathfrak{p}}\Bigr),
\]
其中 $r_0$ 是整除 $\mathfrak{m}$ 的实素个数, $h$ 是 $K$ 的类数 ($C$ 的阶).
\end{theorem}

\begin{proof}
包含映射 $I^{S(\mathfrak{m})}\to I$ 定义了一个同态 $C_{\mathfrak{m}}\to C$. 考虑映射对
\[
K_{\mathfrak{m},1}\xrightarrow{\ f\ }K_{\mathfrak{m}}\xrightarrow{\ g\ }I^{S(\mathfrak{m})}.
\]
由引理 \ref{lem:v.1.1}, $g$ 的核与余核分别为 $U$ 与 $C$. $g\circ f$ 的余核是 $C_{\mathfrak{m}}$ (按定义), 其核是 $K_{\mathfrak{m},1}\cap U=U_{\mathfrak{m},1}$. 最后 $f$ 是单射. 因此该映射对的核--余核序列 (见 II, \ref{lem:ii.A.2}) 为
\[
0\to U_{\mathfrak{m},1}\to U\to K_{\mathfrak{m}}/K_{\mathfrak{m},1}\to C_{\mathfrak{m}}\to C\to0.
\]
下面证明 $K_{\mathfrak{m}}$ 与所给群典范同构. 设 $\mathfrak{p}$ 是整除 $\mathfrak{m}$ 的素. 若 $\mathfrak{p}$ 为实素, 则把 $\alpha\in K_{\mathfrak{m}}$ 映到 $\alpha_{\mathfrak{p}}$ 的符号 (回顾: 实素是一个嵌入 $K\hookrightarrow\R$, 而 $\alpha_{\mathfrak{p}}$ 表示 $\alpha$ 在该嵌入下的像). 若 $\mathfrak{p}$ 为有限素, 即它是 $\mathcal{O}_K$ 中的素理想, 则把 $\alpha\in K_{\mathfrak{m}}$ 映到 $[a][b]^{-1}\in(\mathcal{O}_K/\mathfrak{p}^{m(\mathfrak{p})})^{\times}$, 其中 $a,b$ 如引理 \ref{lem:v.1.5} 所述. 由于 $a$ 与 $b$ 都和 $\mathfrak{p}$ 互素, 它们在 $\mathcal{O}_K/\mathfrak{p}^{m(\mathfrak{p})}$ 中的类 $[a]$ 与 $[b]$ 可逆, 故此映射有意义. 弱逼近定理 (ANT, 7.20) 表明映射 $K_{\mathfrak{m}}\to\prod\{\pm\}\times\prod(\mathcal{O}_K/\mathfrak{p}^{m(\mathfrak{p})})^{\times}$ 是满射, 其核显然是 $K_{\mathfrak{m},1}$.

中国剩余定理表明有环的同构
\[
\mathcal{O}_K/\mathfrak{m}_0\simeq\prod_{\mathfrak{p}|\mathfrak{m}}\mathcal{O}_K/\mathfrak{p}^{m(\mathfrak{p})},
\]
从而有群的同构
\[
(\mathcal{O}_K/\mathfrak{m}_0)^{\times}\simeq\prod(\mathcal{O}_K/\mathfrak{p}^{m(\mathfrak{p})})^{\times}.
\]
这就完成了同构部分的证明.

剩下计算群的阶. 注意 $\mathcal{O}_K/\mathfrak{p}^m$ 是局部环, 极大理想为 $\mathfrak{p}/\mathfrak{p}^m$ (因为其理想对应于 $\mathcal{O}_K$ 中包含 $\mathfrak{p}^m$ 的理想), 故其单位是不属于 $\mathfrak{p}/\mathfrak{p}^m$ 的元素. 滤链
\[
(\mathcal{O}_K/\mathfrak{p}^m)^{\times}\supset(1+\mathfrak{p})/\mathfrak{p}^m\supset\cdots\supset(1+\mathfrak{p}^{m-1})/\mathfrak{p}^m\supset 0
\]
的商群同构于
\[
k^{\times},\,k,\,\ldots,\,k,\qquad k\overset{\mathrm{def}}{=}\mathcal{O}_K/\mathfrak{p},
\]
故 $(\mathcal{O}_K/\mathfrak{p}^m)^{\times}$ 的阶为 $(q-1)q^{m-1}$, $q=(\mathcal{O}_K:\mathfrak{p})\overset{\mathrm{def}}{=}\Nm\mathfrak{p}$. 这表明
\[
(C_{\mathfrak{m}}:1)=(C:1)\cdot(K_{\mathfrak{m}}:K_{\mathfrak{m},1})\cdot(U:U_{\mathfrak{m},1})^{-1}
\]
等于定理陈述中的表达式.
\end{proof}

\begin{example}\label{ex:v.1.8}
\begin{enumerate}[(a)]
\item 若 $\mathfrak{m}=1$, 则 $C_{\mathfrak{m}}=C$.
\item 当 $\mathfrak{m}$ 是所有实素之积时, $C_{\mathfrak{m}}$ 是\textbf{狭窄类群} (narrow class group)\index{leiqun@类群(class group)!xiazhai@狭窄类群(narrow class group)}, 且有正合列
\[
0\to U/U_+\to K^{\times}/K_+\to C_{\mathfrak{m}}\to C\to1,
\]
其中 $K_+$ 是全正元素群 (即在所有实嵌入下均为正), $U_+$ 是全正单位群. 此外 $K^{\times}/K_+\simeq\prod_{\mathfrak{p}\,\mathrm{real}}\{\pm\}$, 故 $C_{\mathfrak{m}}\to C$ 的核是所有可能的符号模去由单位产生的那些.

对 $\Q$, 狭窄类群平凡. 对 $\Q[\sqrt{d}]$, $d>0$, 有两个实素, 且 $U=\{\pm\varepsilon^m\mid m\in\Z\}\approx(\Z/2\Z)\times\Z$, 其中 $\varepsilon$ 是基本单位. 设 $\bar{\varepsilon}$ 为 $\varepsilon$ 的共轭. 则依 $\varepsilon$ 与 $\bar{\varepsilon}$ 符号相同或不同, 有 $h_m=2h$ 或 $h$. 注意 $\Nm(\varepsilon)=+1$ 当符号相同, $-1$ 当符号相反. 对小的 $d$ 值:
\[
\begin{array}{cccc}
d & h & \varepsilon & \Nm(\varepsilon)\\
2 & 1 & 1+\sqrt{2} & -1\\
3 & 1 & 2+\sqrt{3} & 1\\
5 & 1 & (1+\sqrt{5})/2 & -1\\
6 & 1 & 5+2\sqrt{6} & 1
\end{array}
\]
因此 $\Q[\sqrt{3}]$ 与 $\Q[\sqrt{6}]$ 的类数为 $1$ 而狭窄类数为 $2$, 而对 $\Q[\sqrt{2}]$ 与 $\Q[\sqrt{5}]$, 两个类数都是 $1$.
\item 对域 $\Q$ 与模 $(m)$, 该正合列成为
\[
0\to\{\pm1\}\to(\Z/m\Z)^{\times}\to C_{\mathfrak{m}}\to 0.
\]
对模 $\infty(m)$, 正合列成为
\[
0\to\{\pm1\}\to\{\pm\}\times(\Z/m\Z)^{\times}\to C_{\mathfrak{m}}\to 0.
\]
此处 $-1$ 映到 $(-,[-1])$, 而积的子群 $(\Z/m\Z)^{\times}$ 同构地映满商群 $C_{\mathfrak{m}}$.
\end{enumerate}
\end{example}

\subsection*{Frobenius 元}

我们回顾 Frobenius 元的理论 (ANT, 第 8 章). 设 $K$ 为数域, $L$ 为 $K$ 的有限 Galois 扩张, 群为 $G$. 设 $\mathfrak{p}$ 是 $K$ 的一个素理想, $\mathfrak{P}$ 是 $L$ 中位于其上的素理想. 分解群 $D(\mathfrak{P})$ (或 $G(\mathfrak{P})$) 定义为
\[
\{\tau\in G\mid\tau\mathfrak{P}=\mathfrak{P}\}.
\]
等价地, 它是 $G$ 中对 $\mathfrak{P}$-adic 拓扑连续, 从而延拓到完备化 $L_{\mathfrak{P}}$ 的元素之集. 由此得到同构
\[
D(\mathfrak{P})\to\Gal(L_{\mathfrak{P}}/K_{\mathfrak{p}}).
\]
设 $\mathfrak{P}$ 在 $\mathfrak{p}$ 上不分歧. 则 $\Gal(L_{\mathfrak{P}}/K_{\mathfrak{p}})$ 在 $\mathcal{O}_L$ 上的作用诱导同构
\[
\Gal(L_{\mathfrak{P}}/K_{\mathfrak{p}})\to\Gal(l/k),
\]
其中 $l$ 与 $k$ 是剩余类域. 图示如下:
\[
\begin{tikzcd}[column sep=small, row sep=small]
\mathfrak{P} \ar[dd, no head] & L \ar[r, no head] \ar[dd, "f", "D(\mathfrak{P})"'] & L_{\mathfrak{P}} \ar[r, no head] \ar[dd, "f", "D(\mathfrak{P})"'] & l \ar[dd, no head, "D(\mathfrak{P})"'] \\
& & & \\
\mathfrak{P}_D \ar[dd, no head] & L^{D(\mathfrak{P})} \ar[r, no head] \ar[dd, "g"'] & K_{\mathfrak{p}} \ar[r, no head] \ar[ddl, no head] & k \\
& & & \\
\mathfrak{p} & K \ar[r, no head] & \mathfrak{p} &
\end{tikzcd}
\]
群 $\Gal(l/k)$ 是循环群, 有典范生成元, 即 Frobenius 元 $x\mapsto x^q$, 其中 $q$ 是 $k$ 的元素个数. 从而 $D(\mathfrak{P})$ 循环, $D(\mathfrak{P})$ 中对应于 $\Gal(l/k)$ 里 Frobenius 元的生成元称为在 $\mathfrak{P}$ 处的 \textbf{Frobenius 元}\index{fuluobeniusiyuan@Frobenius 元(Frobenius element)} $(\mathfrak{P},L/K)$. 它是 $\Gal(L/K)$ 中唯一满足下列两个条件的元素 $\sigma$:

(a) $\sigma\in D(\mathfrak{P})$, 即 $\sigma\mathfrak{P}=\mathfrak{P}$;

(b) 对所有 $\alpha\in\mathcal{O}_L$, $\sigma\alpha\equiv\alpha^q\bmod\mathfrak{P}$, 其中 $q$ 是剩余类域 $\mathcal{O}_K/\mathfrak{p}$, $\mathfrak{p}=\mathfrak{P}\cap K$ 的元素个数.

现在列举 $(\mathfrak{P},L/K)$ 的基本性质.

\textbf{1.9.} 设 $\tau\mathfrak{P}$ 是整除 $\mathfrak{p}$ 的另一个素. 则 $D(\tau\mathfrak{P})=\tau D(\mathfrak{P})\tau^{-1}$, 且
\[
(\tau\mathfrak{P},L/K)=\tau(\mathfrak{P},L/K)\tau^{-1}.
\]

\begin{proof}
若 $\rho\in D(\mathfrak{P})$, 则
\[
\tau\rho\tau^{-1}(\tau\mathfrak{P})=\tau\rho\mathfrak{P}=\tau\mathfrak{P},
\]
故 $\tau\rho\tau^{-1}\in D(\tau\mathfrak{P})$. 于是 $\tau D(\mathfrak{P})\tau^{-1}\subset D(\tau\mathfrak{P})$, 又因两者阶相同, 必相等.

设 $\alpha\in\mathcal{O}_L$, $\sigma=(\mathfrak{P},L/K)$; 则
\[
\tau\sigma\tau^{-1}(\alpha)=\tau\bigl((\tau^{-1}\alpha)^q+a\bigr),\quad\text{某个 }a\in\mathfrak{P},
\]
而
\[
\tau\bigl((\tau^{-1}\alpha)^q+a\bigr)=\alpha^q+\tau a\equiv\alpha^q\bmod\tau\mathfrak{P}.
\]
\end{proof}

由于 $G$ 在整除 $\mathfrak{p}$ 的素上传递作用, 这表明
\[
\{(\mathfrak{P},L/K)\mid\mathfrak{P}|\mathfrak{p}\}
\]
是 $G$ 中的一个共轭类, 我们记之为 $(\mathfrak{p},L/K)$. 当 $L/K$ 为 Abel 扩张时, $(\mathfrak{p},L/K)$ 只含一个元素, 我们把它视为 $\Gal(L/K)$ 的一个元素 (而非由单个元素构成的集合).

\phantomsection\hypertarget{item:v.1.10}{}\textbf{1.10.} 考虑域塔
\[
\begin{tikzcd}
M \ar[d, no head] & \Omega \ar[d, no head] \\
L \ar[d, no head] & \mathfrak{P} \ar[d, no head] \\
K & \mathfrak{p}
\end{tikzcd}
\]
其中 $\Omega$ 是 $M$ 中位于 $\mathfrak{p}$ 之上的素理想, $\mathfrak{P}=\Omega\cap\mathcal{O}_L$. 设 $\Omega$ 在 $\mathfrak{p}$ 上不分歧; 则
\[
(\Omega,M/L)=(\Omega,M/K)^{f(\mathfrak{P}/\mathfrak{p})}.
\]

\begin{proof}
设 $k(\Omega)\supset k(\mathfrak{P})\supset k(\mathfrak{p})$ 为剩余类域塔. 则 $f(\mathfrak{P}/\mathfrak{p})\overset{\mathrm{def}}{=}[k(\mathfrak{P}):k(\mathfrak{p})]$, 且 $\Gal(k(\Omega)/k(\mathfrak{P}))$ 中的 Frobenius 元是 $\Gal(k(\Omega)/k(\mathfrak{p}))$ 中 Frobenius 元的 $f(\mathfrak{P}/\mathfrak{p})$ 次幂. 其余是直接的.
\end{proof}

\phantomsection\hypertarget{item:v.1.11}{}\textbf{1.11.} 在 \hyperlink{item:v.1.10}{1.10} 中设 $L$ 在 $K$ 上 Galois; 则
\[
(\Omega,M/K)\big|_L=(\mathfrak{P},L/K).
\]

\begin{proof}
显然 $(\Omega,M/K)|_L$ 满足刻画 $(\mathfrak{P},L/K)$ 的条件.
\end{proof}

设 $L_1$ 与 $L_2$ 是含于某个域 $\Omega$ 的 $K$ 的 Galois 扩张, 并令 $M=L_1\cdot L_2$. 则 $M$ 在 $K$ 上 Galois, 且有单同态
\[
\sigma\mapsto(\sigma|L_1,\sigma|L_2)\colon\Gal(M/K)\to\Gal(L_1/K)\times\Gal(L_2/K).
\]

\phantomsection\hypertarget{item:v.1.12}{}\textbf{1.12.} 设 $\mathfrak{Q}$ 是 $\mathcal{O}_M$ 的素理想, $\mathfrak{P}_i=\mathfrak{Q}\cap\mathcal{O}_{L_i}$. 在上述映射下,
\[
(\mathfrak{Q},M/K)\mapsto(\mathfrak{P}_1,L_1/K)\times(\mathfrak{P}_2,L_2/K).
\]

\begin{proof}
应用 \hyperlink{item:v.1.11}{1.11}.
\end{proof}

注意 $\mathfrak{p}$ 在 $L$ 中完全分裂当且仅当 $(\mathfrak{P},L/K)=1$ 对某个 (从而所有) 位于其上的 $\mathfrak{P}$ 成立. 因此在 \hyperlink{item:v.1.12}{1.12} 的情形, $\mathfrak{p}$ 在 $M$ 中完全分裂当且仅当它在 $L_1$ 与 $L_2$ 中都完全分裂.

\section{L-级数与算术级数中素数的密度}\label{sec:V.2}

我们先简要回顾 Dirichlet $L$-级数的初等理论 (例如见 \cite{Serre1970}, 第 VI 章).

设 $m$ 为整数. \textbf{模 $m$ 的 Dirichlet 特征}\index{dirichlettezheng@Dirichlet 特征(Dirichlet character)}是同态 $\chi\colon(\Z/m\Z)^{\times}\to\CC^{\times}$. 因 $(\Z/m\Z)^{\times}$ 有限, $\chi([n])$ 对所有 $n$ 是 $1$ 的根. 模 $m$ 的 Dirichlet 特征可视为与 $m$ 互素的整数集上的乘性函数, 其在 $n$ 处的值只依赖 $n$ 模 $m$. 常把 $\chi$ 延拓到全体整数上, 规定当 $\gcd(m,n)\ne 1$ 时 $\chi(n)=0$. 模 $m$ 的平凡 Dirichlet 特征, 即在与 $m$ 互素的所有整数上取值 $1$, 称为主 Dirichlet 特征 $\chi_0$.

对模 $m$ 的 Dirichlet 特征 $\chi$, 附带一个 Dirichlet 级数
\[
L(s,\chi)=\prod_{p\nmid m}\frac{1}{1-\chi(p)p^{-s}}=\sum_{n>0}\chi(n)/n^s.
\]
两个表达式都对 $\Re(s)>1$ 的复数 $s$ 收敛, 二者的相等正是唯一分解的解析表述. 注意 $L(s,\chi_0)$ 与 Riemann zeta 函数 $\zeta(s)$ 仅相差被积的因子 $\frac{1}{1-p^{-s}}$ (对整除 $m$ 的 $p$).

\begin{theorem}\label{thm:v.2.1}
\begin{enumerate}[(a)]
\item zeta 函数 $\zeta(s)$ 可延拓为半平面 $\Re(s)>0$ 上的亚纯函数, 且
\[
\zeta(s)=\frac{1}{s-1}+\varphi(s),
\]
其中 $\varphi(s)$ 在 $\Re(s)>0$ 上全纯.
\item 若 $\chi\ne\chi_0$, 则 $L(s,\chi)$ 的级数对 $\Re(s)>0$ 收敛, 且 $L(1,\chi)\ne 0$.
\end{enumerate}
\end{theorem}

\begin{proof}
见 \cite{Serre1970}, 第 VI 章, 命题 10, 12, 定理 1.
\end{proof}

对 (a) 中的等式取对数, 利用近似 $-\log(1-x)\sim x$ 可得
\[
\sum 1/p^s\sim\log\frac{1}{s-1}\quad(s\downarrow 1).\erratum{原书此处及下文密度定义、定理 \ref{thm:v.2.2} 的证明中共三处误作 $\log\frac{1}{1-s}$, 依 Milne errata (Yiming Tang) 改正.}
\]
我们的意思是: 商 $\dfrac{\sum 1/p^s}{-\log(s-1)}$ 当 $s$ 从大于 $1$ 的实数一侧趋近 $1$ 时收敛到 $1$. 这一结果使如下定义显得合理: 若
\[
\sum_{p\in T}1/p^s\sim\delta\log\frac{1}{s-1}\quad(s\downarrow 1),
\]
则称素数集 $T$ 具有 \textbf{Dirichlet (或解析) 密度} $\delta$.

定义 $f_{\chi}(s)=\sum_{p\nmid m}\chi(p)/p^s$. 则 (b) 表明, 对 $\chi\ne\chi_0$, $f_{\chi}(s)$ 在 $s=1$ 附近有界. 一个初等论证 (\cite{Serre1970}, 第 VI 章, 引理 9) 表明, 对每个与 $m$ 互素的 $a$,
\[
\sum_{p\equiv a\bmod m}1/p^s=\frac{1}{\varphi(m)}\sum_{\chi}\chi(a)^{-1}f_{\chi}(s),
\]
其中 $\varphi(m)\overset{\mathrm{def}}{=}|(\Z/m\Z)^{\times}|$, 求和遍历模 $m$ 的所有 Dirichlet 特征.

\begin{theorem}\label{thm:v.2.2}
对每个与 $m$ 互素的 $a$, 算术级数
\[
\ldots,\ a-2m,\ a-m,\ a,\ a+m,\ a+2m,\ \ldots
\]
中的素数具有 Dirichlet 密度 $1/\varphi(m)$.
\end{theorem}

\begin{proof}
对 $\chi\ne\chi_0$, $f_{\chi}(s)$ 在 $s=1$ 附近有界, 故
\[
\sum_{p\equiv a\bmod m}1/p^s\sim\frac{1}{\varphi(m)}\chi_0(a)^{-1}f_{\chi_0}(s)\sim\frac{1}{\varphi(m)}\log\frac{1}{s-1}\quad(s\downarrow 1).
\]
\end{proof}

\begin{corollary}\label{cor:v.2.3}
对每个 $m$, 在分圆域 $\Q[\zeta_m]$ 中分裂的素数集具有 Dirichlet 密度 $1/\varphi(m)$.
\end{corollary}

\begin{proof}
素理想 $(p)$ 在 $\Q[\zeta_m]$ 中分裂当且仅当 $p\equiv 1\bmod m$.
\end{proof}

现在说明上述结果如何从 $\Q$ 推广到任意数域. 证明将在第 \ref{ch:VI} 章给出.

设 $K$ 为数域, $\mathfrak{m}$ 为 $K$ 的模. \textbf{模 $\mathfrak{m}$ 的 (Dirichlet 或 Weber) 特征}是同态 $\chi\colon C_{\mathfrak{m}}\to\CC^{\times}$, 其值同样是 $1$ 的根. 换一种说法, 这样的特征是一个乘性函数 $I^S\to\CC^{\times}$, 它在 $i(K_{\mathfrak{m},1})$ 上取零值, 对某个满足 $S(\mathfrak{m})=S$ 的模 $\mathfrak{m}$. 模 $\mathfrak{m}$ 的主特征是只取值 $1$ 的函数 $\chi_0\colon C_{\mathfrak{m}}\to\CC^{\times}$.

对模 $\mathfrak{m}$ 的特征 $\chi$, 附带一个 Dirichlet 级数
\[
\mathcal{L}(s,\chi)=\prod_{\mathfrak{p}\nmid\mathfrak{m}}\frac{1}{1-\chi(\mathfrak{p})\Nm\mathfrak{p}^{-s}}=\sum_{(\mathfrak{a},\mathfrak{m}_0)=\mathcal{O}_K}\chi(\mathfrak{a})/\Nm\mathfrak{a}^s.
\]
乘积遍历与 $\mathfrak{m}_0$ 互素的素理想, 求和遍历与 $\mathfrak{m}_0$ 互素的整理想. 同样, 两个表达式都对 $\Re(s)>1$ 收敛, 其相等正是理想唯一分解的解析表述. $L$-级数 $\mathcal{L}(s,\chi_0)$ 与 Dedekind zeta 函数
\[
\zeta_K(s)\overset{\mathrm{def}}{=}\prod_{\mathfrak{p}}\frac{1}{1-\Nm\mathfrak{p}^{-s}}=\sum_{\mathfrak{a}\subset\mathcal{O}_K}\Nm\mathfrak{a}^{-s}
\]
仅相差有限多个因子.

\begin{theorem}\label{thm:v.2.4}
\begin{enumerate}[(a)]
\item zeta 函数 $\zeta_K(s)$ 可延拓为半平面 $\Re(s)>0$ 上的亚纯函数, 且当 $s\downarrow 1$ 时
\[
\zeta_K(s)\sim\frac{2^{r_1}(2\pi)^{r_2}\,\mathrm{Reg}(K)}{w_K|\Delta_{K/\Q}|^{1/2}}\,h_K\cdot\frac{1}{s-1},
\]
其中 $r_1$ 与 $r_2$ 分别是 $K$ 的实素与复素个数, $\mathrm{Reg}(K)$ 是 $K$ 的 regulator (见 ANT, 第 5 章), $w_K$ 是 $K$ 中 $1$ 的根的个数, $\Delta_{K/\Q}$ 是 $K/\Q$ 的判别式, $h_K$ 是类数.
\item 若 $\chi\ne\chi_0$, 则 $L(s,\chi)$ 的级数对 $\Re(s)>0$ 收敛, 且 $L(1,\chi)\ne 0$.
\end{enumerate}
\end{theorem}

(b) 的证明用到存在定理 (见 \ref{thm:v.3.6}).

同样, 对 (a) 中的等式取对数可得
\[
\sum_{\mathfrak{p}}1/\Nm\mathfrak{p}^s\sim\log\frac{1}{s-1}\quad(s\downarrow 1),
\]
且若
\[
\sum_{\mathfrak{p}\in T}1/\Nm\mathfrak{p}^s\sim\delta\log\frac{1}{s-1}\quad(s\downarrow 1),
\]
则称 $K$ 中素理想集 $T$ 具有 \textbf{Dirichlet (或自然) 密度} $\delta$.

与前一情形类似的论证证明:

\begin{theorem}\label{thm:v.2.5}
对每个与 $\mathfrak{m}_0$ 互素的理想 $\mathfrak{a}$, $\mathcal{O}_K$ 中在 $C_{\mathfrak{m}}$ 里类为 $[\mathfrak{a}]$ 的素理想具有 Dirichlet 密度 $1/h_{\mathfrak{m}}$.
\end{theorem}

定理 \ref{thm:v.2.4} 与 \ref{thm:v.2.5} 证明中的分析与 $K=\Q$ 的情形相同, 但数论部分要困难得多.

\section{理想语言的主要定理}\label{sec:V.3}

\subsection*{Artin 映射}

设 $L/K$ 为有限 Abel 扩张, Galois 群为 $G$. 回顾: 对在 $L$ 中不分歧的素理想 $\mathfrak{p}$, 有唯一确定的 Frobenius 自同构 $\sigma=(\mathfrak{p},L/K)\in\Gal(L/K)$, 由下述条件刻画: 对每个位于 $\mathfrak{p}$ 之上的素理想 $\mathfrak{P}$, $\sigma\mathfrak{P}=\mathfrak{P}$, 且 $\sigma\alpha\equiv\alpha^{\Nm\mathfrak{p}}\bmod\mathfrak{P}$.

对包含所有在 $L$ 中分歧的素的 $K$ 的素的每个有限集 $S$, 有同态
\[
\psi_{L/K}\colon I^S\to\Gal(L/K),\qquad\mathfrak{p}_1^{n_1}\cdots\mathfrak{p}_t^{n_t}\mapsto\prod(\mathfrak{p}_i,L/K)^{n_i},
\]
称为\textbf{整体 Artin 映射} (或互反映射)\index{artinyingshe@Artin 映射(Artin map)!zhengti@整体 Artin 映射(global Artin map)}\index{huxiangyingshe@互反映射(reciprocity map)}.

\begin{example}\label{ex:v.3.1}
设 $K=\Q[\sqrt{m}]$, $m$ 为无平方因子整数. 在 $K$ 中分歧的有限素集 $S$: 当 $m\equiv 1\bmod 4$ 时由整除 $m$ 的素组成, 否则由整除 $m$ 的素连同 $2$ 组成. 把 $\Gal(K/\Q)$ 与 $\{\pm1\}$ 等同. Artin 映射是由
\[
p\mapsto\Bigl(\frac{m}{p}\Bigr)\colon I^S\to\Gal(K/\Q)
\]
确定的同态, 其中 $\bigl(\frac{m}{p}\bigr)$ 是二次剩余 (Legendre) 符号.
\end{example}

\begin{example}\label{ex:v.3.2}
设 $L=\Q[\zeta_m]$, $\zeta_m$ 为本原 $m$ 次单位根. 设 $m$ 为奇数或被 $4$ 整除 (此时在 $L$ 中分歧的素恰为整除 $m$ 的素). 把与 $m$ 互素的整数 $n$ 送到 $L$ 的自同构 $\zeta\mapsto\zeta^n$ 的映射定义了一个同构 $(\Z/m\Z)^{\times}\to\Gal(L/\Q)$. 对不整除 $n$ 的 $p$, $(p,L/K)=[p]$. 若 $r$ 与 $s$ 是与 $m$ 互素的正整数, 则 $r/s$ 定义了一个类 $[r/s]=[r][s]^{-1}\in(\Z/m\Z)^{\times}$, 而 Artin 映射是复合
\[
I^S\xrightarrow{\ (r/s)\mapsto[r/s]\ }(\Z/m\Z)^{\times}\xrightarrow{\ [n]\mapsto(\xi\mapsto\xi^n)\ }\Gal(L/\Q).
\]
\end{example}

回顾 (ANT, p.~68): 对数域的任意有限扩张 $L/K$, 范理想映射 $\Nm_{L/K}\colon I_L\to I_K$ 是从 $L$ 的分式理想群到 $K$ 的相应群的唯一同态, 使得对 $L$ 的任意素理想 $\mathfrak{P}$, $\Nm_{L/K}(\mathfrak{P})=\mathfrak{p}^{f(\mathfrak{P}/\mathfrak{p})}$, 其中 $\mathfrak{p}=\mathfrak{P}\cap\mathcal{O}_K$. 对任意 $\alpha\in L$, $\Nm_{L/K}((\alpha))=(\Nm_{L/K}\alpha)$.

\begin{proposition}\label{prop:v.3.3}
设 $L$ 是 $K$ 的 Abel 扩张, $K'$ 是中间域: $L\supset K'\supset K$. 则下图交换:
\[
\begin{tikzcd}
I_{K'}^S\ar[r, "\psi_{L/K'}"]\ar[d, "\Nm"'] & \Gal(L/K')\ar[d, "\text{限制}\ |\!\ _K"]\\
I_K^S\ar[r, "\psi_{L/K}"'] & \Gal(L/K)
\end{tikzcd}
\]
此处 $S$ 是 $K$ 的素理想的任意有限集, 包含所有在 $L$ 中分歧者, 以及 $K'$ 中位于 $S$ 的素之上的素集.
\end{proposition}

\begin{proof}
设 $\mathfrak{p}'$ 是 $K'$ 的任意素理想, 位于 $K$ 的不在 $S$ 中的素理想 $\mathfrak{p}$ 之上. 则 $\Nm_{K'/K}(\mathfrak{p}')=\mathfrak{p}^{f(\mathfrak{p}'/\mathfrak{p})}$, 须证 $\psi_{L/K'}(\mathfrak{p}')=\psi_{L/K}(\mathfrak{p}^{f(\mathfrak{p}'/\mathfrak{p})})$, 即对每个位于 $\mathfrak{p}$ 之上的 $L$ 的素理想 $\mathfrak{P}$, $(\mathfrak{P},L/K')=(\mathfrak{P},L/K)^{f(\mathfrak{p}'/\mathfrak{p})}$. 但这正是 \hyperlink{item:v.1.10}{1.10} 中证明过的.
\end{proof}

\begin{corollary}\label{cor:v.3.4}
对 $K$ 的每个有限 Abel 扩张 $L$, $\Nm_{L/K}(I_L^S)$ 包含在 $\psi_{L/K}\colon I^S\to\Gal(L/K)$ 的核中.
\end{corollary}

\begin{proof}
在上图中取 $K'=L$.
\end{proof}

于是, 当 $L/K$ 为有限 Abel 扩张时, Artin 映射诱导同态
\[
\psi_{L/K}\colon I_K^S/\Nm(I_L^S)\to\Gal(L/K).
\]
群 $I^S/\Nm(I_L^S)$ 是无限的 (因为有无穷多个素理想不分裂), 故 $\psi_{L/K}$ 不可能单.

\subsection*{整体类域论的主要定理}

设 $S$ 为 $K$ 的素的有限集. 称同态 $\psi\colon I^S\to G$ \textbf{容许一个模}, 若存在模 $\mathfrak{m}$, $S(\mathfrak{m})\supset S$, 使得 $\psi(i(K_{\mathfrak{m},1}))=0$. 换言之, $\psi$ 容许一个模当且仅当它经过某个满足 $S(\mathfrak{m})\supset S$ 的 $C_{\mathfrak{m}}$ 分解.

\begin{theorem}[互反律]\label{thm:v.3.5}
设 $L$ 是 $K$ 的有限 Abel 扩张, $S$ 是 $K$ 中在 $L$ 分歧的素集. 则 Artin 映射 $\psi\colon I^S\to\Gal(L/K)$ 容许一个满足 $S(\mathfrak{m})=S$ 的模 $\mathfrak{m}$, 且它定义同构
\[
I_K^{S(\mathfrak{m})}/i(K_{\mathfrak{m},1})\cdot\Nm(I_L^{S(\mathfrak{m})})\to\Gal(L/K).
\]
\end{theorem}

如定理陈述中的模称为 $L$ 的\textbf{定义模} (defining modulus).

注意该定理并未蕴含 $K$ 有哪怕一个非平凡 Abel 扩张. 记 $I_K^{\mathfrak{m}}$ 为 $K$ 中 $S(\mathfrak{m})$-理想群, $I_L^{\mathfrak{m}}$ 为 $L$ 中 $S(\mathfrak{m})'$-理想群, 其中 $S(\mathfrak{m})'$ 包含位于 $S(\mathfrak{m})$ 的素之上的 $L$ 的素. 称 $I_K^{\mathfrak{m}}$ 的子群 $H$ 为\textbf{模 $\mathfrak{m}$ 的同余子群}, 若
\[
I_K^{\mathfrak{m}}\supset H\supset i(K_{\mathfrak{m},1}).
\]

\begin{theorem}[存在定理]\label{thm:v.3.6}
对模 $\mathfrak{m}$ 的每个同余子群 $H$, 存在有限 Abel 扩张 $L/K$, 以 $\mathfrak{m}$ 为模 (即在不整除 $\mathfrak{m}$ 的素处不分歧), 使得
\[
H=i(K_{\mathfrak{m},1})\cdot\Nm_{L/K}(I_L^{\mathfrak{m}}).\erratum{原书未指明 $L/K$ 以 $\mathfrak{m}$ 为模; 若无此条件, 存在定理并不能 (立即) 推出下文关于 Artin 映射诱导同构的断言. 依 Milne errata (Zhuo Zhang) 补正.}
\]
\end{theorem}

注意, 对如定理中的 $H$ 与 $L$, Artin 映射 $\psi_{L/K}$ 诱导同构
\[
I^{S(\mathfrak{m})}/H\to\Gal(L/K).
\]

特别地, 对每个模 $\mathfrak{m}$ 存在一个域 $L_{\mathfrak{m}}$, 称为\textbf{模 $\mathfrak{m}$ 的射线类域}\index{shexianleiyu@射线类域(ray class field)}, 使得 Artin 映射定义同构 $C_{\mathfrak{m}}\to\Gal(L_{\mathfrak{m}}/K)$. 对域 $L\subset L_{\mathfrak{m}}$, 令
\[
\Nm(C_{L,\mathfrak{m}})=i(K_{\mathfrak{m},1})\cdot\Nm(I_L^{\mathfrak{m}})\ \mathrm{mod}\ i(K_{\mathfrak{m},1}).
\]

\begin{corollary}\label{cor:v.3.7}
固定模 $\mathfrak{m}$. 则映射 $L\mapsto\Nm(C_{L,\mathfrak{m}})$ 是从含于 $L_{\mathfrak{m}}$ 的 $K$ 的 Abel 扩张之集到 $C_{\mathfrak{m}}$ 的子群之集的双射. 此外,
\begin{align*}
L_1\subset L_2&\iff\Nm(C_{L_1,\mathfrak{m}})\supset\Nm(C_{L_2,\mathfrak{m}});\\
\Nm(C_{L_1\cdot L_2,\mathfrak{m}})&=\Nm(C_{L_1,\mathfrak{m}})\cap\Nm(C_{L_2,\mathfrak{m}});\\
\Nm(C_{L_1\cap L_2,\mathfrak{m}})&=\Nm(C_{L_1,\mathfrak{m}})\cdot\Nm(C_{L_2,\mathfrak{m}}).
\end{align*}
\end{corollary}

\begin{remark}\label{rem:v.3.8}
设 $L/K$ 是 Galois 群为 $G$ 的 Abel 扩张. 据互反律, 存在模 $\mathfrak{m}$, 其支撑为在 $L$ 分歧的素集, 使 Artin 映射 $\psi_{L/K}\colon I^{S(\mathfrak{m})}\to G$ 在 $i(K_{\mathfrak{m},1})$ 上取值 $1$. 考虑定理 \ref{thm:v.1.7} 中的映射
\[
(\mathcal{O}_K/\mathfrak{p}^{m(\mathfrak{p})})^{\times}\hookrightarrow K_{\mathfrak{m}}/K_{\mathfrak{m},1}\xrightarrow{\ i\ }C_{\mathfrak{m}}\xrightarrow{\ \psi_{L/K}\ }G.
\]
显然存在最小整数 $f(\mathfrak{p})\le m(\mathfrak{p})$ 使该映射经过 $(\mathcal{O}_K/\mathfrak{p}^{f(\mathfrak{p})})^{\times}$ 分解. 模 $\mathfrak{f}(L/K)=\mathfrak{m}_\infty\prod\mathfrak{p}^{f(\mathfrak{p})}$ 便是使 $\psi_{L/K}$ 经过 $C_{\mathfrak{f}}$ 分解的最小模, 称为 $L/K$ 的\textbf{导子}\index{daozi@导子(conductor)}\footnote{德语为 F\"uhrer\,---\,在 1930 年代的德国, 公开谈论类域论可能有风险.}. 导子 $\mathfrak{f}(L/K)$ 恰被在 $L$ 中分歧的素整除.

含 $K$ 的射线类域 $L_{\mathfrak{m}}$ 的子域, 恰是导子整除 $\mathfrak{m}$ 的那些域.\footnote{设 $L$ 是导子整除 $\mathfrak{m}$ 的 Abel 扩张. 则在 $L$ 中分裂的素也在 $L_{\mathfrak{m}}$ 中分裂. 再由 Chebotarev 密度定理即得 $L\subset L_{\mathfrak{m}}$ (见下文 \ref{thm:v.3.25}).} $K$ 的每个 Abel 扩张都含于某个 $L_{\mathfrak{m}}$.
\end{remark}

下文第五节将用 idèle 重述定理\ref{thm:v.3.5} 与 \ref{thm:v.3.6}, 并在第 \ref{ch:VII} 章证明重述后的定理.

如下文所讨论的, Artin 映射是满射有一个相当简单的解析证明. 因此证明互反律的难点在于证明 Artin 映射容许一个导子, 且
\[
\bigl(I^{S(\mathfrak{m})}:i(K_{\mathfrak{m},1})\cdot\Nm(I_L^{S(\mathfrak{m})})\bigr)\le[L:K].
\]
为证明存在定理, 须对每个模构造射线类域. 遗憾的是, 我们不知道如何直接构造射线类域; 实际做法是构造足够多的扩张, 迫使定理成立.

\begin{example}\label{ex:v.3.9}
模 $\mathfrak{m}=1$ 的射线类群是理想类群, 相应的射线类域是 \textbf{Hilbert 类域}\index{hilbertleiyu@Hilbert 类域(Hilbert class field)}; 它是 $K$ 的极大 Abel 扩张, 在包括实素的所有素处不分歧 (即实素保持实). 例如, $\Q$ 的 Hilbert 类域是 $\Q$ 自身 (因 $\Q$ 类数为 $1$). $\Q[\sqrt{-5}]$ 的 Hilbert 类域是 $\Q[\sqrt{-1},\sqrt{5}]$\,---\,$2$ 与 $5$ 都在 $\Q[\sqrt{-5}]$ 中分歧, 但在 $\Q[\sqrt{-1}]$ 中只有 $2$ 分歧, 在 $\Q[\sqrt{5}]$ 中只有 $5$ 分歧, 由此推出 $\Q[\sqrt{-5}]$ 中整除 $2$ 与 $5$ 的素不在 $\Q[\sqrt{-1},\sqrt{5}]$ 中分歧.
\end{example}

\begin{example}\label{ex:v.3.10}
设 $m$ 为奇数或被 $4$ 整除的正整数. 模 $(m)$ 的射线类域是 $\Q[\zeta_m+\bar{\zeta}_m]$, 模 $\infty(m)$ 的射线类域是 $\Q[\zeta_m]$.\footnote{这由 \ref{ex:v.1.8}(c) 中对射线类群的描述以及 ANT, 8.18 中对 Frobenius 元的描述推出.} 于是互反律蕴含 Kronecker--Weber 定理: $\Q$ 的每个 Abel 扩张的导子整除某个 $\infty(m)$, 从而含于某个分圆域.
\end{example}

\begin{example}\label{ex:v.3.11}
设 $d$ 为无平方因子整数. 我们通过求使 $\Q[\sqrt{d}]\subset\Q[\zeta_m]$ 的最小整数 $m$ 来计算 $K=\Q[\sqrt{d}]$ 在 $\Q$ 上的导子.

先考虑奇素数 $p$. 此时 $\Gal(\Q[\zeta_p]/\Q)\simeq(\Z/p\Z)^{\times}$ 是 $p-1$ 阶循环群, 故有唯一的 $2$ 阶商群. 因此 $\Q[\zeta_p]$ 含唯一的二次域, 由于它只能在 $p$ 处分歧, 必为 $\Q[\sqrt{p^*}]$, 其中 $p^*=(-1)^{\frac{p-1}{2}}p$ (符号的选取使 $p^*\equiv 1\bmod 4$).

其次注意 $\zeta_8=(1+i)/\sqrt{2}$, 故 $\zeta_8+\bar{\zeta}_8=\sqrt{2}$. 因此 $\Q[\sqrt{2}]\subset\Q[\zeta_8]$ (事实上, $\Q[\sqrt{2}]$ 是 $\Q[\zeta_8]$ 的极大实子域, 且 $\Q[\zeta_8]=\Q[i,\sqrt{2}]$).

设 $n$ 为整除 $d$ 的奇素数之积 (故 $d=\pm n$ 或 $\pm 2n$). 我断言
\[
\Q[\sqrt{d}]\subset\Q[\zeta_n]\quad\text{若 } d\equiv 1\bmod 4,
\]
\[
\Q[\sqrt{d}]\subset\Q[\zeta_{4n}]\quad\text{若 } d\equiv 3\bmod 4,
\]
\[
\Q[\sqrt{d}]\subset\Q[\zeta_{8n}]\quad\text{若 } d\equiv 2\bmod 4,
\]
且每种情形这都是包含 $\Q[\sqrt{d}]$ 的最小分圆域. 例如, 注意 $d=p_1\cdots p_r$, $d\equiv 1\bmod 4$, 蕴含 $d=p_1^*\cdots p_r^*$, 故 $\Q[\sqrt{d}]\subset\Q[\zeta_n]$. 又注意若 $d$ 为偶数, 则 $\Q[\sqrt{d}]$ 不含于 $\Q[\zeta_{4n}]$, 否则 $\Q[\zeta_{4n}]$ 将包含 $i,\sqrt{d},\sqrt{d/2}$, 从而包含 $i,\sqrt{2}$ 与 $\zeta_8$.

结论: $\Q[\sqrt{d}]$ 的导子是 $|\Delta_{K/\Q}|$ 或 $\infty|\Delta_{K/\Q}|$, 视 $d>0$ 或 $d<0$ 而定, 此处 $\Delta_{K/\Q}$ 是 $K/\Q$ 的判别式.
\end{example}

\begin{remark}\label{rem:v.3.12}
对有限 Abel 扩张 $L/K$ 及满足 $S(\mathfrak{m})$ 等于 $K$ 中在 $L$ 分歧的素集的模 $\mathfrak{m}$, 令
\[
T(L/K,\mathfrak{m})=i(K_{\mathfrak{m},1})\cdot\Nm_{L/K}(I_L^{S(\mathfrak{m})})\subset I_K.
\]
Takagi 证明了: 当 $\mathfrak{m}$ 充分可除时 (事实上, 对所有被 $L/K$ 的导子 $\mathfrak{f}$ 整除的 $\mathfrak{m}$), 群 $T(L/K,\mathfrak{m})$ 不依赖于 $\mathfrak{m}$, 且 $I_K^{S(\mathfrak{m})}/T(L/K,\mathfrak{m})\approx\Gal(L/K)$. 因此当 $\mathfrak{f}|\mathfrak{m}$ 时, $T(L/K,\mathfrak{m})$ 常被称为 $L/K$ 的 \textbf{Takagi 群}. Artin 证明了 Takagi 群正是映射 $\mathfrak{a}\mapsto(\mathfrak{a},L/K)\colon I_K^{S(\mathfrak{m})}\to\Gal(L/K)$ 的核.
\end{remark}

\begin{exercise}\label{exc:v.3.13}
利用二次互反律求使 $i(\Q_{m,1})$ 含于 Artin 映射之核的最小 $m$, 以此计算 $\Q[\sqrt{d}]/\Q$ 的导子 (对照习题 \ref{ex:0.11}).
\end{exercise}

对应于同余子群 $H$ 的域 $L$ 称为 $H$ 的\textbf{类域}, 这也是本学科名称的由来. 注意: 对不整除 $L/K$ 的导子的 $K$ 的素 $\mathfrak{p}$, 位于 $\mathfrak{p}$ 之上的素的剩余类次数 $f(\mathfrak{P}/\mathfrak{p})$ 等于 $\mathfrak{p}$ 在 $I^{\mathfrak{m}}/H$ 中的阶 (因为这是 $\Gal(k(\mathfrak{P})/k(\mathfrak{p}))$ 中 Frobenius 元的阶). 于是我们得到了用 $K$ 的理想结构对 $K$ 的 Abel 扩张的分类, 并且对每个 Abel 扩张都知道 $K$ 中素的分解规律.

\begin{exercise}\label{exc:v.3.14}
验证下表的最后一行:
\begin{center}
\begin{tabular}{|c|c|c|c|c|c|}
\hline
判别式 & $-15$ & $-20$ & $-23$ & $-24$ & $-31$\\
\hline
类数 & $2$ & $2$ & $3$ & $2$ & $3$\\
\hline
Hilbert 类域 & $X^2+3$ & $X^2+1$ & $X^3-X-1$ & $X^2+3$ & $X^3+X-1$\\
\hline
\end{tabular}
\end{center}
第一行列出前五个类数不为 $1$ 的虚二次域的判别式, 第二行是它们的类数, 最后一行是 Hilbert 类域生成元的极小多项式. (注意对全虚域, 类数与狭窄类数相等.)

[例如, 判别式为 $-24$ 的二次域是 $\Q[\sqrt{-6}]$, 其 Hilbert 类域是 $L=\Q[\sqrt{-6},\sqrt{-3}]=\Q[\sqrt{2},\sqrt{-3}]$. 于是 $L=\Q[\alpha]$, $\alpha=\sqrt{2}+\sqrt{-3}$, 其极小多项式为 $X^4+2X^2+5$. 然而 $\Z[\alpha]$ 并不等于 $L$ 的整数环; 对照 \mo{122765}.]
\end{exercise}

\begin{exercise}\label{exc:v.3.15}
本习题解释当我们忽略 $K$ 的素理想的有限集 $S$ 时会发生什么. 设 $\mathfrak{m}$ 是 $K$ 的满足 $S(\mathfrak{m})\cap S=\emptyset$ 的模, $H$ 是 $I^{S\cup S(\mathfrak{m})}$ 的包含 $i(K_{\mathfrak{m},1})$ 的子群. 称 $K$ 的扩张 $L$ 为 $H$ 的 \textbf{$S$-类域}, 若

(a) $L$ 是 $K$ 的有限 Abel 扩张, 且 $S$ 中的素理想在 $L$ 中完全分裂;

(b) $m(\mathfrak{p})=0\Rightarrow\mathfrak{p}$ 不在 $L$ 中分歧;

(c) 不在 $S\cup S(\mathfrak{m})$ 中且在 $L$ 中分裂的素理想恰是 $H$ 中的那些.

证明: 对如上每个群 $H$ 都存在 $S$-类域 $L$, 它是唯一的, 且 $I^{S\cup S(\mathfrak{m})}/H\simeq\Gal(L/K)$; 此外, 每个满足 (a) 的域 $L$ 都是某个 $H$ 的 $S$-类域.

提示: 证明 $I^{S\cup S(\mathfrak{m})}/i(K_{\mathfrak{m},1})\simeq I^{S(\mathfrak{m})}/\langle S\rangle\cdot i(K_{\mathfrak{m},1})$, 其中 $\langle S\rangle$ 是由 $S$ 中的素生成的 $I^{S(\mathfrak{m})}$ 的子群.
\end{exercise}

\subsection*{范限制定理}

在对 $K$ 的 Abel 扩张的分类中, 我们对 $L$ 附上群 $H=i(K_{\mathfrak{m},1})\cdot\Nm(I_L^{S(\mathfrak{m})})$, 其中 $\mathfrak{m}$ 是充分大到成为定义模的模 (此时 $(I^{S(\mathfrak{m})}:H)=[L:K]$). 人们或许希望类似的办法对非 Abel 扩张也有效, 但下面的定理表明并非如此.

\begin{theorem}[范限制定理]\label{thm:v.3.16}
设 $L$ 是 $K$ 的有限扩张, $L'/K$ 是 $L/K$ 的极大 Abel 子扩张. 对 $L'/K$ 的每个定义模 $\mathfrak{m}$,
\[
i(K_{\mathfrak{m},1})\cdot\Nm_{L/K}(I_L^{S(\mathfrak{m})})=i(K_{\mathfrak{m},1})\cdot\Nm_{L'/K}(I_{L'}^{S(\mathfrak{m})}).
\]
\end{theorem}

例如, 若 $L$ 是 $K$ 的在 $K$ 上非 Galois 的三次扩张, 则 $L'=K$, 故
\[
i(K_{\mathfrak{m},1})\cdot\Nm_{L/K}(I_L^{S(\mathfrak{m})})=I_K^{S(\mathfrak{m})}.
\]

范限制定理表明: 对非 Abel 扩张 $L/K$, $\mathrm{Spl}(L/K)$ 不能用同余条件描述. 该定理的证明见例如 \cite{GrantLeitzel1969}.

\subsection*{主理想定理}

\begin{quote}
\itshape ... \`a ma honte, je ne suis par arriv\'e \`a retrouver ce \guillemotleft corollaire\guillemotright\ que tous les id\'eaux de $K$ deviennent principaux dans la plus grande extension ab\'elienne non ramifi\'ee dans le fini. Si \c ca s'explique en deux mots, je t'en serais fort reconnaissant.
\end{quote}
\hspace*{\fill}Grothendieck, 致 Serre 的信, 1956.9.19\footnote{\guillemotleft\,\ldots\ 令我羞愧的是, 我一直没能重新找出这个\guillemotleft 推论\guillemotright: $K$ 的所有理想在有限素处不分歧的极大 Abel 扩张中都变为主理想. 若能用两句话解释, 我将不胜感激.\guillemotright\ Serre 回信道: \guillemotleft 随信附上一篇关于 \guillemotleft Hauptidealsatz\guillemotright\ 的小文章, 解释该定理如何化归为一个 (实际上非常神秘的) 群论定理. 事实上, 这正是 Artin 本人在其论文 (Abh.~Hamburg, 大约第 7--10 卷) 中给出的化归; 如果你能找到该定理的一个漂亮的上同调证明, 那就再好不过了, 但迄今为止人人都被卡住了.\guillemotright\ 另见 Serre 信件的尾注 (\cite{GrothendieckSerre2001}).}

下面的定理由 Hilbert 在 1900 年前后猜想.

\begin{theorem}[主理想定理]\label{thm:v.3.17}
$K$ 的每个理想在 $K$ 的 Hilbert 类域中都变为主理想.
\end{theorem}

我解释一下证明的想法. 回顾: 对群 $G$, $G$ 的\textbf{换位子群} (或\textbf{导群}) $G'$ 是由换位子 $ghg^{-1}h^{-1}$, $g,h\in G$, 生成的子群. 商群 $G^{\mathrm{ab}}=G/G'$ 是 Abel 群, 且是 $G$ 的极大 Abel 商. 若 $L$ 是 $K$ 的 Galois 扩张, Galois 群为 $G$, 则 $L^{G'}$ 是 $K$ 的 Galois 群为 $G^{\mathrm{ab}}$ 的 Abel 扩张, 且是 $L$ 中包含 $K$ 的极大 Abel 扩张.

设有域
\[
L\supset K'\supset K,
\]
其中 $L$ 在 $K$ 上 Galois (不必 Abel). 对 $K$ 的素的每个有限集 $S$, $\mathfrak{a}\mapsto\mathfrak{a}\mathcal{O}_{K'}$ 是同态 $I_K^S\to I_{K'}^S$. 考虑:
\[
\begin{tikzcd}
I_K^S\ar[r, "\psi_{L/K}"]\ar[d, "\text{典范}"'] & \Gal(L/K)^{\mathrm{ab}}\ar[d, "?"]\\
I_{K'}^S\ar[r, "\psi_{L/K'}"'] & \Gal(L/K')^{\mathrm{ab}}
\end{tikzcd}
\]
使图交换的映射 ``?'' 是什么?

在描述它之前, 需要解释群论中的一个构造. 设 $H$ 是群 $G$ 的有限指数子群, 把 $G$ 写成陪集的不交并
\[
G=Hg_1\cup Hg_2\cup\cdots\cup Hg_n.
\]
对 $g\in G$, 当 $g\in Hg_i$ 时令 $\varphi(g)=g_i$, 并定义
\[
V(g)=\prod_{i=1}^{n}g_i g\,\varphi(g_i g)^{-1}\ \mathrm{mod}\ H',
\]
其中 $H'$ 是 $H$ 的换位子群.

\begin{proposition}\label{prop:v.3.18}
映射 $g\mapsto V(g)$ 是同态 $G\to H/H'$, 且不依赖陪集代表元 $g_i$ 的选取.
\end{proposition}

\begin{proof}
验证是直接的\,---\,例如见 \cite{Hall1959}, 14.2.1.
\end{proof}

于是, 只要有了群 $G$ 及其有限指数子群 $H$, 就有确定的同态
\[
V\colon G^{\mathrm{ab}}\to H^{\mathrm{ab}},
\]
称为 \textbf{Verlagerung} (或\textbf{转移}) 映射\index{zhuanyiying@转移映射(transfer homomorphism)}\index{verlagerung@Verlagerung(Verlagerung)}. 粗略地说, 我们想定义一个从 $G$ 到 $H$ 的\guillemotleft 范\guillemotright, 由于群不交换这行不通, 但过渡到相应的 Abel 群后就行得通了.

在上图的情形, $\Gal(L/K')$ 是 $\Gal(L/K)$ 的子群, 从而 Verlagerung 是同态
\[
V\colon\Gal(L/K)^{\mathrm{ab}}\to\Gal(L/K')^{\mathrm{ab}}.
\]
Emil Artin 证明了这正是使上图交换的映射 (对照 II, \ref{prop:ii.3.2}b). 考虑域
\[
K''\supset K'\supset K,
\]
其中 $K'$ 是 $K$ 的 Hilbert 类域, $K''$ 是 $K'$ 的 Hilbert 类域. 则

(a) $K''$ 在 $K$ 上正规 (因为 $K''$ 的任何共轭仍是 $K'$ 的 Abel 非分歧扩张, 从而含于 $K''$);

(b) $K'$ 是 $K''$ 中包含 $K$ 的极大 Abel 扩张 (因为 $K''$ 中包含 $K$ 的每个 Abel 扩张在 $K$ 上非分歧, 从而含于 $K'$).

由 (b) 得 $\Gal(K''/K)^{\mathrm{ab}}=\Gal(K'/K)$. 因此取 $L=K''$ 时, 前图成为
\[
\begin{tikzcd}
C_K\ar[r, "\simeq"]\ar[d, "\text{典范}"'] & \Gal(K'/K)\ar[d, "\mathrm{Ver}"]\\
C_{K'}\ar[r, "\simeq"'] & \Gal(K''/K')
\end{tikzcd}
\]
其中 $C_K$ 与 $C_{K'}$ 是 $K$ 与 $K'$ 的类群. 设 $G=\Gal(K''/K)$, $H=\Gal(K''/K')$. 由 (b), $H$ 是 $G$ 的换位子群. 下面的定理 (由 Emil Artin 猜想) 表明此情形下 $V$ 为零, 从而典范映射 $C_K\to C_{K'}$ 为零, 即 $K$ 的每个理想在 $K'$ 中变为主理想.

\begin{theorem}\label{thm:v.3.19}
设 $G$ 是有限群, $H$ 是其换位子群; 则
\[
V\colon G^{\mathrm{ab}}\to H^{\mathrm{ab}}
\]
为零.
\end{theorem}

\begin{proof}
这是群论的定理, 也称为主理想定理. 由 Furtw\"angler 于 1930 年证明. 简单证明见: Witt, Proc.~International Congress of Mathematicians, Amsterdam, 1954, Vol.~2, pp.~71--73. 另见 Zassenhaus, Theory of Groups, V, 定理 12, 或 \cite{ArtinTate1961}, 第 XII 章.
\end{proof}

\begin{remark}\label{rem:v.3.20}
事实上容易看出: 存在 $K$ 的一个次数整除类数 $h$ 的扩张 $L$, 使 $K$ 的每个理想在 $L$ 中变为主理想: 把 $K$ 的类群写成循环群直和; 对每个直和项取生成元 $a_i$, 设 $a_i$ 在类群中的阶为 $h_i$; 写 $a_i^{h_i}=(a_i)$, 并定义 $L$ 为对每个 $i$ 向 $K$ 添加 $a_i$ 的 $h_i$ 次根所得的域.

然而如此构造的域通常并非 $K$ 的 Hilbert 类域\,---\,它甚至不必在 $K$ 上 Galois. 可能存在次数 $<h$ 的 $K$ 的域, 使 $K$ 的每个理想在其中变为主理想.
\end{remark}

\begin{remark}\label{rem:v.3.21}
注意主理想定理说的是: $K$ 中元素的\guillemotleft 理想\guillemotright\ 分解在 Hilbert 类域 $L$ 中成为实际的分解: 设 $a\in K$, $(a)=\mathfrak{p}_1^{r_1}\cdots\mathfrak{p}_s^{r_s}$; 则 $\mathfrak{p}_i\mathcal{O}_L=(\pi_i)$ 对某个 $\pi_i\in\mathcal{O}_L$ 成立, 于是在 $\mathcal{O}_L$ 中 $a=u\pi_1^{r_1}\cdots\pi_s^{r_s}$.
\end{remark}

\begin{remark}\label{rem:v.3.22}
主理想定理当然并不蕴含 $K$ 的 Hilbert 类域 $K'$ 的每个理想都是主理想, 因为并非 $K'$ 的每个理想都属于同态 $I_K\to I_{K'}$ 的像. 可以继续作 $K'$ 的 Hilbert 类域 $K''$, 等等, 得到塔
\[
K\subset K'\subset K''\subset\cdots\subset K^{(n)}\subset\cdots
\]
其中 $K^{(n+1)}$ 是 $K^{(n)}$ 的 Hilbert 类域. 对 $K^{\mathrm{al}}$ 的每个自同构 $\tau$, $\tau K^{(n+1)}$ 显然是 $\tau K^{(n)}$ 的极大 Abel 非分歧扩张, 这蕴含 $\tau K'=K'$, $\tau K''=K''$, 等等. 特别地, $K^{(n)}$ 在 $K$ 上 Galois. \textbf{类域塔问题} (Hasse 于 1925 年提出) 问: 该塔是否总是有限, 从而终止于一个类数为一的域? Golod 与 Shafarevich 于 1964 年证明答案是否定的 (见 \cite{Roquette1967}). 例如 $\Q[\sqrt{-2\cdot 3\cdot 5\cdot 7\cdot 11\cdot 13}]$ 有无限类域塔. 事实上, 只要 $d$ 的素因子多于 $8$ 个, $\Q[\sqrt{d}]$ 就有无限类域塔.
\end{remark}

\subsection*{Chebotarev 密度定理}

设 $L$ 是 $K$ 的 Galois 扩张, Galois 群为 $G$. 回顾: 对在 $L$ 中不分歧的 $K$ 的每个素理想 $\mathfrak{p}$,
\[
(\mathfrak{p},L/K)\overset{\mathrm{def}}{=}\{(\mathfrak{P},L/K)\mid\mathfrak{P}|\mathfrak{p}\}
\]
是 $G$ 中的一个共轭类.

\begin{theorem}[Chebotarev 密度定理]\label{thm:v.3.23}
设 $L/K$ 是 Galois 群为 $G$ 的数域有限扩张, $C$ 是 $G$ 中的一个共轭类. 则满足 $(\mathfrak{p},L/K)=C$ 的 $K$ 的素理想之集在 $K$ 的所有素理想之集中的密度为 $|C|/|G|$. 特别地, 若 $G$ 为 Abel 群, 则对固定的 $\tau\in G$, 满足 $(\mathfrak{p},L/K)=\tau$ 的 $K$ 的素理想 $\mathfrak{p}$ 之集的密度为 $1/(G:1)$.
\end{theorem}

\begin{proof}
对 Abel 扩张 $L/K$, 这由定理 \ref{thm:v.2.5} 与定理 \ref{thm:v.3.5} 推出: 后者说映射 $\mathfrak{p}\mapsto(\mathfrak{p},L/K)$ 对某个模 $\mathfrak{m}$ 诱导满同态 $C_{\mathfrak{m}}\to\Gal(L/K)$, 前者说素理想在 $C_{\mathfrak{m}}$ 的类中均匀分布. 非 Abel 情形由 Abel 情形经一个巧妙的论证导出\,---\,见第 \ref{ch:VIII} 章.
\end{proof}

\begin{corollary}\label{cor:v.3.24}
若多项式 $f(X)\in K[X]$ 对 $K$ 中除有限多个素理想外的所有素理想 $\mathfrak{p}$ 都模 $\mathfrak{p}$ 分解为一次因子, 则它在 $K[X]$ 中分裂.
\end{corollary}

\begin{proof}
对 $f(X)$ 的分裂域应用定理.
\end{proof}

对数域的有限扩张 $L/K$ 及 $K$ 的素的有限集 $S$, 设 $\mathrm{Spl}_S(L/K)$ 为 $K$ 中不在 $S$ 内且在 $L$ 中分裂的素之集.

\begin{theorem}\label{thm:v.3.25}
若 $L$ 与 $M$ 都是 $K$ 的 Galois 扩张, 则
\[
L\subset M\iff\mathrm{Spl}_S(L/K)\supset\mathrm{Spl}_S(M/K).
\]
从而
\[
L=M\iff\mathrm{Spl}_S(L/K)=\mathrm{Spl}_S(M/K).
\]
\end{theorem}

\begin{proof}
作为 (\hyperlink{item:v.1.12}{1.12}) 的推论,
\[
\mathrm{Spl}_S(LM/K)=\mathrm{Spl}_S(L/K)\cap\mathrm{Spl}_S(M/K).
\]
于是
\[
\mathrm{Spl}_S(L/K)\supset\mathrm{Spl}_S(M/K)\Rightarrow\mathrm{Spl}_S(LM/K)=\mathrm{Spl}_S(M/K),
\]
由 Chebotarev 密度定理, 这蕴含
\[
[LM:K]=[M:K],
\]
故 $L\subset M$. 反向蕴含是显然的.
\end{proof}

\begin{remark}\label{rem:v.3.26}
\begin{enumerate}[(a)]
\item 没有 Galois 假设时定理 \ref{thm:v.3.25} 不成立 (见 Cassels and Fr\"ohlich 1967, p.~363).
\item 在定理 \ref{thm:v.3.25} 的陈述中, $S$ 可换成密度为 $0$ 的任意素集.
\item 设 $f(X)$ 是 $K[X]$ 中的不可约多项式. 若 $f(X)$ 对几乎所有素理想 $\mathfrak{p}$ 都模 $\mathfrak{p}$ 有根, 则 $f(X)$ 在 $K$ 中有根 (Cassels and Fr\"ohlich 1967, p.~363, 6.2).
\item Chebotarev 用 Dirichlet (解析) 密度证明了其定理. Artin 指出它对自然密度 (见下文) 也应成立, 而事实上确实如此.
\item 给定数域间的有限 Galois 扩张 $K/\Q$ 及共轭类 $C$, Chebotarev 定理说存在 $K$ 的素理想 $\mathfrak{p}$, 它在 $\Q$ 上不分歧, 使 $(\mathfrak{p},L/K)=C$. 有趣的是 $\Nm\mathfrak{p}$ 需要有多大.  关于此问题的近作, 见 Ahn, Jeoung-Hwan; Kwon, Soun-Hi.~An explicit upper bound for the least prime ideal in the Chebotarev density theorem.~Ann.~Inst.~Fourier (Grenoble) 69 (2019), no.~3, 1411--1458.
\end{enumerate}
\end{remark}

\subsection*{导子--判别式公式}

两个 Dirichlet 特征 $\chi\colon I^S\to\CC^{\times}$ 与 $\chi'\colon I^{S'}\to\CC^{\times}$ 称为\textbf{相合的} (cotrained), 若它们在某个 $I^{S''}$, $S''\supset S\cup S'$, 上一致. 这是一个等价关系. 每个等价类中存在唯一的 $\chi$, 其 $S$ 最小\,---\,这样的 $\chi$ 称为\textbf{本原的} (primitive).

设 $\chi_1$ 是与 $\chi$ 等价的本原特征. 使 $\chi_1$ 在 $i(K_{\mathfrak{m},1})$ 上取零值的最小模 $\mathfrak{m}$ 称为 $\chi$ 的\textbf{导子} $\mathfrak{f}(\chi)$. 令 $\mathfrak{f}(\chi)=\mathfrak{f}_\infty(\chi)\mathfrak{f}_0(\chi)$, 其中 $\mathfrak{f}_\infty$ 与 $\mathfrak{f}_0$ 分别只被无限素与有限素整除.

\begin{theorem}[F\"uhrerdiskriminantenproduktformel]\label{thm:v.3.27}
对 Galois 群为 $G$ 的数域的每个有限 Abel 扩张 $L/K$,
\begin{align*}
\mathrm{disc}(L/K)&=\prod_{\chi\in G^{\vee}}\mathfrak{f}_0(\chi\circ\psi_{L/K}),\qquad G^{\vee}\overset{\mathrm{def}}{=}\mathrm{Hom}(G,\CC^{\times});\\
\mathfrak{f}(L/K)&=\mathrm{lcm}_{\chi}\bigl(\mathfrak{f}(\chi\circ\psi_{L/K})\bigr).
\end{align*}
\end{theorem}

显然 $\bigcap\mathrm{Ker}(\chi\colon G\to\CC^{\times})=0$, 由此得第二个等式. 第一个等式的证明从略\,---\,它实际上是关于局部域的陈述.

\begin{remark}\label{rem:v.3.28}
设 $H$ 是特征 $\chi$ 的核. 则 $\mathfrak{f}(\chi)=\mathfrak{f}(L^H/K)$. 例如, 若 $\chi$ 单, 则 $\mathfrak{f}(\chi)=\mathfrak{f}(L/K)$.
\end{remark}

\begin{example}\label{ex:v.3.29}
设 $L=\Q[\sqrt{d}]$. 则 $G\simeq\Z/2\Z$, 只有两个特征 $G\to\CC^{\times}$: 平凡特征 $\chi_0$ 与单射特征 $\chi_1$. 因此定理说 $\Delta_{K/\Q}=\mathfrak{f}_0(\chi_1)=\mathfrak{f}_0(K/\Q)$, 与例 \ref{ex:v.3.11} 中所证一致.
\end{example}

\begin{example}\label{ex:v.3.30}
设 $L=\Q[\zeta_p]$, $p$ 为奇素数. 则 $G\simeq(\Z/p\Z)^{\times}$ 是 $p-1$ 阶循环群, 故有 $p-2$ 个非平凡特征. 若 $\chi$ 非平凡, 则 $\mathfrak{f}(\chi)|\infty(p)$, 但 $\mathfrak{f}(\chi)=1$ 或 $\infty$ 不可能 (因为那将蕴含 $\chi$ 平凡). 因此当 $\chi$ 非平凡时 $\mathfrak{f}(\chi)=(p)$ 或 $\infty(p)$, 于是导子--判别式公式 (正确地) 给出 $\Delta_{L/\Q}=\pm p^{p-2}$.
\end{example}

\begin{exercise}\label{exc:v.3.31}
对扩张 $\Q[\zeta_{p^2}]/\Q$ 验证导子--判别式公式.
\end{exercise}

\begin{remark}\label{rem:v.3.32}
设 $L/K$ 是 Galois 群为 $G$ (不必 Abel) 的数域有限扩张. 对 $G$ 在有限维向量空间 $V$ 上的表示 $\rho\colon G\to\mathrm{GL}(V)$, Artin 附带一个 Dirichlet $L$-级数 $L(s,\rho)$ (见引言). 这些 Artin $L$-级数的解析性质至今仍未完全弄清.

当 $G$ 交换时, Artin 映射 $I^S\to G$ 把 $G$ 的特征与 Dirichlet/Weber 特征等同, 从而把 Artin $L$-级数与 Dirichlet $L$-级数等同. 这正是 Artin 寻求该映射的动机 (而不像今天看起来那么自然的那样, 是为了在早已知道抽象同构的群 $C_{\mathfrak{m}}$ 与 $G$ 之间构造典范同构).

Langlands 纲领的一部分正是 Dirichlet $L$-级数与 Abel Artin $L$-级数之间这一对应的巨大推广.
\end{remark}

\subsection*{互反律与幂互反律}

设 $K$ 含本原 $n$ 次单位根, $a\in K$. 若 $\sqrt[n]{a}$ 是 $X^n-a$ 的一个根, 则其余根形如 $\zeta\sqrt[n]{a}$, $\zeta$ 为 $n$ 次单位根. 因此 $L\overset{\mathrm{def}}{=}K[\sqrt[n]{a}]$ 在 $K$ 上 Galois, 且 $\sigma\sqrt[n]{a}=\zeta\sqrt[n]{a}$ 对某个 $n$ 次单位根 $\zeta$ 成立.

若 $\mathfrak{p}$ 是 $K$ 的与 $n$ 和 $a$ 互素的素理想, 则 $\mathfrak{p}$ 在 $L$ 中不分歧, 且可由公式
\[
(\mathfrak{p},L/K)(\sqrt[n]{a})=\Bigl(\frac{a}{\mathfrak{p}}\Bigr)_{\!n}\sqrt[n]{a}
\]
定义一个 $n$ 次单位根 $\bigl(\frac{a}{\mathfrak{p}}\bigr)_n$.

可以证明
\[
\Bigl(\frac{a}{\mathfrak{p}}\Bigr)_{\!n}=1\iff a\ \text{模}\ \mathfrak{p}\ \text{是}\ n\ \text{次幂},
\]
故 $\bigl(\frac{a}{\mathfrak{p}}\bigr)_n$ 推广了二次剩余符号. 因此 $\bigl(\frac{a}{\mathfrak{p}}\bigr)_n$ 称为\textbf{幂剩余符号}\index{mishengyufuhao@幂剩余符号(power residue symbol)}. Artin 互反律蕴含这些符号的所有已知的互反律, 因而正如 Artin 指出的, 可视为它们到没有单位根的域的推广. 我们将在第 \ref{ch:VIII} 章解释这一点.

\subsection*{一个初等的未解决问题}

设 $K$ 为数域, $S$ 为 $K$ 的素理想的非空有限集, $p$ 为不被 $S$ 中任何素整除的素数. 是否存在域列 $\ldots,L_n,L_{n+1},\ldots$ 使得

(a) $L_n$ 在 $S$ 之外的素处不分歧;

(b) $p^n|[L_n:K]$?

一个关键情形 (答案未知) 是 $K=\Q$, $S=\{l\}$.

更明确地 (也稍难些): 固定 $\Q$ 中互异的素数 $p\ne l$. 是否存在首一不可约多项式列 $f_n(X)\in\Z[X]$ 使得

(a) $\mathrm{disc}(f_n(X))$ 不被 $l$ 以外的任何素数整除;

(b) $p^n|\deg f_n(X)$?

见 \cite{Milne2006}, pp.~48--49, 了解这个问题出现于何处. 它当然非常困难.\footnote{设 $K$ 为数域, $S$ 为 $K$ 的素集, $K^S$ 为 $K$ 在 $S$ 外不分歧的极大扩张. 设 $P$ 是使 $\ell^\infty\mid [K_S:K]$ 的素数 $\ell$ 之集. 现已知: 只要 $S$ 包含 $K$ 中位于 $\Q$ 的至少两个素数之上的全部素理想, $P$ 就包含所有素数. 见 Chenevier and Clozel, JAMS 22 (2009), 467--519, Cor.~5.2.}

\subsection*{显式整体类域论: Kronecker 的青春之梦与 Hilbert 第十二问题}

与局部类域论不同, 整体类域论 (一般而言) 并不给出数域 $K$ 的 Abel 扩张的显式构造.

Gauss 知道 $\Q$ 的分圆扩张是 Abel 的. 1840 年代末, Kronecker 萌生了一个想法: 分圆域及其子域穷尽了 $\Q$ 的 Abel 扩张; 而且虚二次数域 $E$ 的每个 Abel 扩张都含于向 $E$ 添加单位根与模函数 $j$ 的某些特殊值所得的扩张中. 许多年后他称这一想法为\guillemotleft 我青春时期最珍爱的梦\guillemotright\ (mein liebster Jugendtraum).

更精确地说, Kronecker 之梦\footnote{关于 Kronecker 的想法及其工作的细致梳理, 见 \cite{Schappacher1998}.} 是: $\Q$ 的每个 Abel 扩张都含于向 $\Q$ 添加函数 $e^{2\pi iz}$ 对所有 $z\in\Q^{\times}$ 的值所得的域; 虚二次域 $K$ 的每个 Abel 扩张都含于向 $K$ 添加函数 $j(z)$ 对所有 $z\in K^{\times}$ 的值所得的域.

后来 Hilbert 把它列入其著名问题的第十二个: 对任意数域 $K$, 找出对 $K$ 起着指数函数对 $\Q$、模函数 $j$ 对虚二次域所起作用的函数.\footnote{见 Mathematical Developments arising from Hilbert's Problems, Proc.~Symp.~Pure Math.~XXVIII, Part 1, AMS, 1976, pp.~18--20.} Kronecker 之梦的第一部分\,---\,$\Q$ 的每个 Abel 扩张都是分圆扩张的子域\,---\,由 Weber (1886, 1899, 1907, 1911) 与 Hilbert (1896) 证明.

上述\guillemotleft Kronecker 之梦\guillemotright\ 的表述并不完全正确. 设 $K$ 为虚二次域. 可写 $\mathcal{O}_K=\Z+\Z\tau$, $\Im(\tau)>0$. 已知 $K[j(\tau)]$ 是 $K$ 的 Hilbert 类域. 现向 $K$ 添加所有单位根及所有值 $j(\tau)$, $\tau\in K$, $\Im(\tau)>0$. 所得的域 $K'$ 在 $K$ 上 Abel, 且 $[K^{\mathrm{ab}}:K']$ 是若干 $2$ 阶群之积. 要得到整个 $K^{\mathrm{ab}}$, 还须添加其他椭圆函数的特殊值. 这些结论由 Weber (1908) 与 Feuter (1914) 部分证明, 由 Takagi (1920) 完全证明.

从现代观点看, 椭圆模函数的特殊值与带复乘的椭圆曲线的算术相关联, 正是关于后者的结果使人得以证明前者生成虚二次域的 Abel 扩张.

自 1950 年代末 Taniyama、Shimura 与 Weil 的工作开始, 椭圆曲线与椭圆模曲线的理论已被成功地推广到高维. 在该理论中, 由虚二次域给出复乘的椭圆曲线被由 CM 域给出复乘的 Abel 簇所取代\,---\,CM 域即全实域 $F$ 的二次全虚扩张 $K$, 而椭圆模函数被自守函数所取代.

在哲学层面, 人们预期: 除 $\Q$ 外, 不能通过添加自守函数的特殊值得到全实域的 Abel 扩张. 然而, 粗略地说, 已知对 CM 域 $K$ 确能得到与该限制相容的最大可能的 Abel 扩张.

更精确地, 设 $K$ 为 CM 域, $F$ 为 $K$ 的极大全实子域. 则 $G\overset{\mathrm{def}}{=}\Gal(\Q^{\mathrm{al}}/K)$ 是 $G'\overset{\mathrm{def}}{=}\Gal(\Q^{\mathrm{al}}/F)$ 中指数为 $2$ 的子群, 相应的 Verlagerung 是同态 $V\colon G'^{\mathrm{ab}}\to G^{\mathrm{ab}}$. 此情形下 $V$ 有非常简单的描述.

\begin{theorem}\label{thm:v.3.33}
设 $K$ 为 CM 域, $F$ 为 $K$ 的全实子域, $[K:F]=2$. 设 $H$ 是 Verlagerung 映射
\[
\Gal(F^{\mathrm{al}}/F)^{\mathrm{ab}}\to\Gal(K^{\mathrm{al}}/K)^{\mathrm{ab}}
\]
的像. 则向 $K$ 添加所有定义在有理权 Shimura 簇的典范模型上的自守函数的特殊值所得的 $K$ 的扩张是 $(K^{\mathrm{ab}})^H\cdot\Q^{\mathrm{ab}}$.
\end{theorem}

\begin{proof}
见 \cite{Wei1993}, \cite{Wei1994}.
\end{proof}

\subsection*{注}

同余群与 $K$ 的 Abel 扩张之间的联系在 Artin 定义其映射之前就已被知晓. 它在约 1870--1920 年间缓慢浮现, 主要贡献者是 Kronecker, Weber, Hilbert 与 Takagi. Chebotarev 于 (1926) 年证明了其定理 (Frobenius 更早证明了一个较不精确的结果), 而 Artin 于 1927 年定义了他的映射并证明它给出同构. (更早, 人们已知 $I^m/H\approx\Gal(L/K)$, 但并无典范同构.) 证明这些形式上纯代数的主要定理竟需要分析\,---\,以 Chebotarev (或 Frobenius) 定理的形态出现\,---\,被视为一个缺陷; 经过大量努力, Chevalley 于 1940 年成功地给出了主要定理的纯代数证明. (难点在于证明: 若 $L/K$ 是素次数 $p$ 的数域 Abel 扩张, 则 $K$ 中至少有一个素在 $L$ 中既不分裂也不分歧.\footnote{类域论对 $k$ 有限的 $k(X)$ 的有限扩张也成立; 局部类域论对剩余类域 $k$ 仅为拟有限的局部域, 即满足 $\Gal(k^{\mathrm{al}}/k)\simeq\hat{\Z}$ 的完全域, 也成立. 当人们试图对 $k$ 拟有限的 $k(X)$ 的有限扩张建立整体类域论时, 发现一切都行得通, 只是对某些 $k$ 确实存在每个素都分裂的非平凡 Abel 扩张. 见本书 Arithmetic Duality Theorems, 第 I 章附录.}) 他还引入了 idèle, 使直接对无限扩张陈述类域论成为可能. 群上同调 (至少 2-上循环等) 从 1920 年代起就在类域论中被隐含地使用, 但在 1950 年代被 Nakayama, Hochschild 与 Tate 系统地使用. 1951/52 年, Artin 与 Tate 在一个影响深远的研究班中给出了类域论的纯代数且极富上同调色彩的处理. 此后, 我们对局部类域论的理解有了重要改进 (主要归功于 Lubin 与 Tate). 非 Abel 类域论是 Langlands 纲领的一部分, 该纲领是一系列相互交织的庞大猜想, 目前已取得一些进展, 尤其是在局部情形与函数域情形. 对更一般的域 ($\Q$ 或 $\F_p$ 上有限超越次数的域), Bloch, Kato, Saito 等人建立了相当令人满意的 Abel 类域论. 它使用了代数 $K$-理论 (综述见 \cite{Raskind1995}, 近期结果见 \cite{Wiesend2007}).

\section{Idèle}\label{sec:V.4}

定理\ref{thm:v.3.5} 与 \ref{thm:v.3.6} 表明: 对任意数域 $K$, 存在典范同构 $\varprojlim_{\mathfrak{m}}C_{\mathfrak{m}}\to\Gal(K^{\mathrm{ab}}/K)$. 与其直接研究 $\varprojlim_{\mathfrak{m}}C_{\mathfrak{m}}$, 不如引入另一个以它为商的群更为自然\,---\,这就是 idèle 类群.

\phantomsection\label{pg:orig169}%
\subsection*{拓扑群}

带拓扑的群 $G$ 称为\textbf{拓扑群}, 若映射
\[
g,g'\mapsto gg'\colon G\times G\to G,\qquad g\mapsto g^{-1}\colon G\to G
\]
连续. 平移映射
\[
g\mapsto ag\colon G\to G
\]
此时是同胚.

一般地, 要在一个集合上确定拓扑, 须给出每点的基本邻域组, 即该点的一组邻域, 使每个邻域都包含组中的一个. 由于平移是同胚, 拓扑群的拓扑由 $1$ 的基本邻域组确定.

我们将用到拓扑群的若干一般性事实, 它们可在许多书中找到. 幸运的是, 我们只需要十分初等的内容.

设 $(X_i)_i$ 是一族 (可以无限) 拓扑空间. $\prod X_i$ 上的\textbf{乘积拓扑}是使形如 $\prod U_i$ ($U_i$ 在 $X_i$ 中开, 且除有限多个 $i$ 外 $U_i=X_i$) 的集合构成基的拓扑. Tychonoff 定理说紧空间的乘积是紧的. 然而, 局部紧空间的无穷乘积一般不再局部紧: 若 $V_i$ 是 $x_i$ 在 $X_i$ 中的紧邻域, 则 $\prod V_i$ 是紧的, 但除非除有限多个 $i$ 外 $V_i=X_i$, 它不是 $(x_i)$ 的邻域.

\subsection*{Idèle}

今后常用 $v$ 记 $K$ 的一个素. 则:

$|\cdot|_v$ = $v$ 的正规化绝对值 (使乘积公式成立的那个);

$K_v$ = $K$ 在 $v$ 处的完备化;

$\mathfrak{p}_v$ = $\mathcal{O}_K$ 中相应的素理想 (当 $v$ 有限时);

$\mathcal{O}_v$ = $K_v$ 的整数环;

$U_v=\mathcal{O}_v^{\times}$;

$\hat{\mathfrak{p}}_v$ = $\mathfrak{p}_v$ 的完备化 = $\mathcal{O}_v$ 的极大理想.

回顾: 对所有 $v$, $K_v$ 局部紧\,---\,事实上 $\mathcal{O}_v$ 是 $0$ 的紧邻域. 类似地 $K_v^{\times}$ 局部紧; 事实上
\[
1+\hat{\mathfrak{p}}_v\supset 1+\hat{\mathfrak{p}}_v^2\supset 1+\hat{\mathfrak{p}}_v^3\supset\cdots
\]
是由开紧子群构成的 $1$ 的基本邻域组.

我们想把所有的群 $K_v^{\times}$ 合成一个大的拓扑空间, 但 $\prod K_v^{\times}$ 不是局部紧的. 于是定义 \textbf{idèle} 群\index{idele@idèle(idèle)}为\footnote{idèle 的其他记号: $\mathbb{A}_K^{\times}$, $\mathrm{GL}_1(\mathbb{A}_K)$.}
\[
\mathbb{I}_K=\Bigl\{(a_v)\in\prod K_v^{\times}\ \Big|\ a_v\in\mathcal{O}_v^{\times}\ \text{对除有限多个外的所有}\ v\Bigr\}.
\]
下文中不带下标的 $\mathbb{I}$ 表示 $\mathbb{I}_K$.

对包含所有无限素的素的每个有限集 $S$, 令
\[
\mathbb{I}_S=\prod_{v\in S}K_v^{\times}\times\prod_{v\notin S}\mathcal{O}_v^{\times},
\]
带乘积拓扑. 第一个因子是局部紧空间的有限乘积, 故局部紧; 第二个因子是紧空间的乘积, 故紧 (由 Tychonoff). 因此 $\mathbb{I}_S$ 局部紧. 注意
\[
\mathbb{I}=\bigcup\mathbb{I}_S.
\]
我们想在 $\mathbb{I}$ 上赋予拓扑, 使每个 $\mathbb{I}_S$ 在 $\mathbb{I}$ 中开且继承乘积拓扑. 做法: 规定开集基由形如 $\prod_v V_v$ 的集合组成, 其中 $V_v$ 在 $K_v^{\times}$ 中开, 且对\textbf{几乎所有}\footnote{此处\guillemotleft 几乎所有\guillemotright\ 意为\guillemotleft 除可能有限多个外的所有\guillemotright.} $v$ 有 $V_v=\mathcal{O}_v^{\times}$. 两个这种形状的集合的交仍包含一个这种形状的集合, 故它们确构成某拓扑的基. 显然该拓扑具有我们想要的性质, 而且它赋予 $\mathbb{I}$ 拓扑群的结构. 下列集合构成 $1$ 的基本邻域组: 对每个包含 $S_\infty$ 的有限素集 $S$ 及 $\varepsilon>0$, 定义
\[
U(S,\varepsilon)=\{(a_v)\mid |a_v-1|_v<\varepsilon,\ v\in S,\quad |a_v|_v=1,\ \text{所有 }v\notin S\}.
\]

\textbf{4.1.} 存在典范的满同态
\[
(a_v)\mapsto\prod_{v\ \mathrm{finite}}\mathfrak{p}_v^{\mathrm{ord}_{\mathfrak{p}}(a_v)}\colon\mathbb{I}_K\to I_K,
\]
其核为 $\mathbb{I}_{S_\infty}$.

可以把 idèle 想成理想的加厚: 它包含无限素的因子, 也包含有限素处的单位. 注意 $\mathbb{I}/\mathbb{I}_{S_\infty}$ 是可数多个 $\Z$ 的带离散拓扑的直和, 而 $\prod K_v^{\times}/\mathbb{I}_{S_\infty}$ 是可数多个 $\Z$ 的直积, 后者本身不可数.

\textbf{4.2.} 存在典范的单射 (对角) 同态
\[
a\mapsto(a,a,a,\ldots)\colon K^{\times}\to\mathbb{I}_K.
\]
我断言其像是离散的. 由于是群, 只须证 $1\in K^{\times}$ 在诱导拓扑下是开的. 取 $U=U(S,\varepsilon)$, $S$ 为任意包含 $S_\infty$ 的有限集, $1>\varepsilon>0$. 对每个 $a\in K^{\times}\cap U$,
\[
\begin{cases}
|a-1|_v<\varepsilon & \text{对所有 } v\in S,\\
|a|_v=1 & \text{对所有 } v\notin S.
\end{cases}
\]
第二个条件蕴含
\[
|a-1|_v\le\max\bigl(|a|_v,|-1|_v\bigr)\le 1.
\]
因此若 $a\in K^{\times}\cap U$, 则 $\prod_v|a-1|_v<\varepsilon^{|S|}<1$, 与乘积公式矛盾, 除非 $a=1$.

商群 $C=\mathbb{I}/K^{\times}$ 称为 $K$ 的 \textbf{idèle 类群}\index{ideleleiqun@idèle 类群(idèle class group)}. 它映满 $K$ 的理想类群. 它不紧 (见下文 (\hyperlink{item:v.4.4}{4.4})).

\textbf{4.3.} 存在典范的单同态
\[
a\mapsto(1,\ldots,1,a,1,\ldots,1)\colon K_v^{\times}\to\mathbb{I}_K
\]
($a$ 在第 $v$ 个位置). $K_v^{\times}$ 上诱导的拓扑是其自然拓扑, 因为
\[
U(S,\varepsilon)\cap K_v^{\times}=
\begin{cases}
\{a\mid|a-1|_v<\varepsilon\} & v\in S,\\
\{a\mid|a|_v=1\} & v\notin S,
\end{cases}
\]
且这些集合构成 $K_v^{\times}$ 中 $1$ 的基本邻域组.

\phantomsection\hypertarget{item:v.4.4}{}\textbf{4.4.} 存在典范的满同态
\[
\mathbf{a}=(a_v)\mapsto c(\mathbf{a})=\prod|a_v|_v\colon\mathbb{I}\to\R_{>0}.
\]
$\mathbf{a}$ 的像称为 $\mathbf{a}$ 的\textbf{容量} (content). 定义
\[
\mathbb{I}^1=\mathrm{Ker}(c)=\{\mathbf{a}\in\mathbb{I}\mid c(\mathbf{a})=1\}.
\]
注意由乘积公式, $K^{\times}\subset\mathbb{I}^1$. 商群 $\mathbb{I}/K^{\times}$ 不可能紧, 因为它满射到 $\R_{>0}$; 但可以证明 $\mathbb{I}^1/K^{\times}$ 是紧的.

\setcounter{theorem}{4} % 4.1--4.4 为叙述性编号条目, 不占定理计数器; 调至 4 使下一条显示为 4.5
\begin{aside}\label{asi:v.4.5}
按同样方式仅用有限素定义的 $\mathbb{I}_f$ 称为\textbf{有限 idèle} 群. 我们有
\[
\prod_{v\ \mathrm{finite}}\mathcal{O}_v^{\times}\subset\mathbb{I}_f\subset\prod_{v\ \mathrm{finite}}K_v^{\times}.
\]
子群 $\prod\mathcal{O}_v^{\times}$ 在 $\mathbb{I}_f$ 中开且紧, 且 $\mathbb{I}_f/\prod\mathcal{O}_v^{\times}=I$ ($K$ 的理想群).

仍有 $K^{\times}$ 到 $\mathbb{I}_f$ 的对角嵌入, 但此时 $K^{\times}$ 上的诱导拓扑可描述如下: $U_K\overset{\mathrm{def}}{=}\mathcal{O}_K^{\times}$ 是开的, 而 $1$ 的基本邻域组由 $U_K$ 的有限指数子群构成 (非平凡定理).\footnote{取自 \sx{140729}. 第一个断言: 有限 idèle 中由 $\prod\mathcal{O}_v^{\times}$ 给出的子集按限制直积拓扑的定义是开的, 而 $K^{\times}$ 中的元素属于该子集当且仅当它在每个完备化中都是单位, 即当且仅当它属于 $\mathcal{O}_K^{\times}$. 第二个断言的证明只是第一个论证的推广: 有限 idèle 中 $1$ 的开邻域基可以这样给出\,---\,在有限多个素 $v$ 处任取 $K_v^{\times}$ 的开子群, 在其余 $v$ 处取 $\mathcal{O}_v^{\times}$. 若只想要开集基, 不妨只考虑非常小的那些, 故可设在所选的有限多个 $v$ 处取 $\mathcal{O}_v^{\times}$ 的小开子集. 注意 $\mathcal{O}_v^{\times}$ 的子群都是有限指数的, 且由素元的幂生成. 把这样的邻域与 $K^{\times}$ 的对角像相交, 得到的是 $\mathcal{O}_K^{\times}$ 中落在所选的那些 $\mathcal{O}_v^{\times}$ 的开子群里的元素构成的子群, 其指数恰为诸指数之积, 故为有限指数子群; 反之, $\mathcal{O}_K^{\times}$ 的任何有限指数子群由其在各完备化中的像确定.} 特别地, $K^{\times}$ 是 $\mathbb{I}_f$ 的离散子群 $\iff U_K$ 有限 $\iff K=\Q$ 或虚二次域.
\end{aside}

\begin{NOTES}
adèle 拓扑并不在 idèle 上诱导 idèle 拓扑 (因为 $\prod_{v|\infty}K_v^{\times}\times\prod_{v\,\mathrm{finite}}\mathcal{O}_v^{\times}$ 在 adèle 中不开). 作为练习, 证明 idèle 在诱导拓扑下甚至不是拓扑群. (先考虑 $\Q$ 的情形, 忘掉无限素. (a) 确保理解 adèle 中 $0$ 的标准邻域基. (b) 把 (a) 中的邻域平移加 $1$, 与 idèle 相交, 得到 adèle 拓扑下 idèle 中 $1$ 的邻域基. (c) 验证 (b) 中某个邻域在映射 $x\mapsto x^{-1}$ 下的原像不开.)
\end{NOTES}

\begin{NOTES}
设 $(V_v,U_v)_v$ 是一族拓扑空间 $V_v$, 每个带有子空间 $U_v$. 该族的\textbf{限制直积} $\prod(V_v,U_v)$ 是 $\prod V_v$ 中由满足\guillemotleft 对几乎所有 $v$, $a_v\in U_v$\guillemotright\ 的族 $(a_v)$ 组成的子集, 其拓扑使 $\prod U_v$ 开且带乘积拓扑. 限制直积有时记作 $\prod' V_v$. 于是 $\mathbb{A}_K=\prod_v(K_v,U_v)$, 其中 $v$ 有限时 $U_v=\mathcal{O}_v$, 否则 $U_v=K_v$. 类似地 $\mathbb{A}_K^{\times}=\prod_v(K_v^{\times},U_v^{\times})$, 其中 $v$ 有限时 $U_v=\mathcal{O}_v^{\times}$, 否则 $U_v=K_v^{\times}$.
\end{NOTES}

\begin{NOTES}
再举些例子. 设 $K$ 是类数为 $1$ 的域. 设 $x=(x_v)$ 是一个 idèle. 则 $x$ 的理想是主理想, 设等于 $(a)$. 此时 $x/a\in\prod_{\mathrm{finite}}\mathcal{O}_v^{\times}\times K_\infty^{\times}$. 故 $\mathbb{A}_K^{\times}=K^{\times}\cdot\prod_{\mathrm{finite}}\mathcal{O}_v^{\times}\times K_\infty^{\times}$. 然而 $K^{\times}\cap(\prod_{\mathrm{finite}}\mathcal{O}_v^{\times}\times K_\infty^{\times})=U_K$.
\end{NOTES}

\begin{NOTES}
取自 \sx{80763}. 普林斯顿资格考试的问题: 关于 adèle 与 idèle 的主要紧性定理给出了哪些结果? 答: 理想类群的有限性与单位定理 (用经典方法证明) 可证范数一 idèle 类群的紧性; 而用测度论方法证明的范数一 idèle 类群的紧性又可反推理想类群的有限性与单位定理. (见 \cite{Cassels1967}.) 另一方面, 紧性结果确实可推广到可除代数, 那里没有明显的理想类群等概念. (见 Weil 的 Basic Number Theory.)
\end{NOTES}

\subsection*{把射线类群实现为 $\mathbb{I}$ 的商}

我们已经看到类群 $C_K=I/i(K^{\times})$ 可以实现为 $\mathbb{I}$ 的商. 现在要对 $C_{\mathfrak{m}}$ 证明同样的事.

设 $\mathfrak{m}$ 为模. 对 $\mathfrak{p}|\mathfrak{m}$, 令
\[
W_{\mathfrak{m}}(\mathfrak{p})=
\begin{cases}
\R_{>0} & \mathfrak{p}\ \text{实},\\
1+\hat{\mathfrak{p}}^{m(\mathfrak{p})} & \mathfrak{p}\ \text{有限}.
\end{cases}
\]
于是在每种情形, $W_{\mathfrak{m}}(\mathfrak{p})$ 都是 $K_{\mathfrak{p}}^{\times}$ 中 $1$ 的邻域.

定义 $\mathbb{I}_{\mathfrak{m}}$ 为满足 $a_{\mathfrak{p}}\in W_{\mathfrak{m}}(\mathfrak{p})$ (对所有 $\mathfrak{p}|\mathfrak{m}$) 的 idèle $(a_{\mathfrak{p}})_{\mathfrak{p}}$ 之集:
\[
\mathbb{I}_{\mathfrak{m}}=\Bigl(\prod_{\mathfrak{p}\nmid\mathfrak{m}}K_{\mathfrak{p}}^{\times}\times\prod_{\mathfrak{p}|\mathfrak{m}}W_{\mathfrak{m}}(\mathfrak{p})\Bigr)\cap\mathbb{I}.
\]
换言之, $\mathbb{I}_{\mathfrak{m}}$ 由 $K$ 的素标号的族 $(a_{\mathfrak{p}})_{\mathfrak{p}}$ 中满足
\[
\begin{cases}
a_{\mathfrak{p}}\in K_{\mathfrak{p}}^{\times} & \text{对所有 }\mathfrak{p},\\
a_{\mathfrak{p}}\in\mathcal{O}_{\mathfrak{p}}^{\times} & \text{对几乎所有 }\mathfrak{p},\\
a_{\mathfrak{p}}\in W_{\mathfrak{m}}(\mathfrak{p}) & \text{对所有 }\mathfrak{p}|\mathfrak{m}
\end{cases}
\]
者组成. 定义 $W_{\mathfrak{m}}$ 为 $\mathbb{I}_{\mathfrak{m}}$ 中使 $a_{\mathfrak{p}}$ 在所有不整除 $\mathfrak{m}$ 的有限 $\mathfrak{p}$ 处为单位的 idèle $(a_{\mathfrak{p}})_{\mathfrak{p}}$ 之集:
\[
W_{\mathfrak{m}}=\prod_{\substack{\mathfrak{p}\nmid\mathfrak{m}\\ \mathfrak{p}\ \mathrm{infinite}}}K_{\mathfrak{p}}^{\times}\times\prod_{\mathfrak{p}|\mathfrak{m}}W_{\mathfrak{m}}(\mathfrak{p})\times\prod_{\substack{\mathfrak{p}\nmid\mathfrak{m}\\ \mathfrak{p}\ \mathrm{finite}}}U_{\mathfrak{p}}.
\]
换言之, $W_{\mathfrak{m}}$ 由满足
\[
\begin{cases}
a_{\mathfrak{p}}\in K_{\mathfrak{p}}^{\times} & \text{对所有无限 }\mathfrak{p},\\
a_{\mathfrak{p}}\in\mathcal{O}_{\mathfrak{p}}^{\times} & \text{对所有有限 }\mathfrak{p},\\
a_{\mathfrak{p}}\in W_{\mathfrak{m}}(\mathfrak{p}) & \text{对所有 }\mathfrak{p}|\mathfrak{m}
\end{cases}
\]
的族 $(a_{\mathfrak{p}})_{\mathfrak{p}}$ 组成. 注意
\[
K_{\mathfrak{m},1}=K^{\times}\cap\prod_{\mathfrak{p}|\mathfrak{m}}W_{\mathfrak{m}}(\mathfrak{p})\qquad(\text{相交在 }\prod_{\mathfrak{p}|\mathfrak{m}}K_{\mathfrak{p}}^{\times}\ \text{内}),
\]
且
\[
K_{\mathfrak{m},1}=K^{\times}\cap\mathbb{I}_{\mathfrak{m}}\qquad(\text{相交在 }\mathbb{I}\ \text{内}).
\]

\begin{proposition}\label{prop:v.4.6}
设 $\mathfrak{m}$ 是 $K$ 的模.

(a) 映射 $\mathrm{id}\colon\mathbb{I}_{\mathfrak{m}}\to I^{S(\mathfrak{m})}$ 定义同构
\[
\mathbb{I}_{\mathfrak{m}}/K_{\mathfrak{m},1}\cdot W_{\mathfrak{m}}\xrightarrow{\ \simeq\ }C_{\mathfrak{m}}.
\]

(b) 包含 $\mathbb{I}_{\mathfrak{m}}\hookrightarrow\mathbb{I}$ 定义同构
\[
\mathbb{I}_{\mathfrak{m}}/K_{\mathfrak{m},1}\to\mathbb{I}/K^{\times}.
\]
\end{proposition}

\begin{proof}
(a) 考虑映射对
\[
K_{\mathfrak{m},1}\to\mathbb{I}_{\mathfrak{m}}\xrightarrow{\ \mathrm{id}\ }I^{S(\mathfrak{m})}.
\]
第一个映射单, 第二个满且核为 $W_{\mathfrak{m}}$, 故该映射对的核--余核序列 (II, \ref{lem:ii.A.2}) 为
\[
W_{\mathfrak{m}}\to\mathbb{I}_{\mathfrak{m}}/K_{\mathfrak{m},1}\to C_{\mathfrak{m}}\to 1.
\]
这证明了命题的 (a).

(b) $\mathbb{I}_{\mathfrak{m}}\to\mathbb{I}/K^{\times}$ 的核是 $K^{\times}\cap\mathbb{I}_{\mathfrak{m}}$ (相交在 $\mathbb{I}$ 内), 如刚才所见它是 $K_{\mathfrak{m},1}$. 于是包含映射定义单射
\[
\mathbb{I}_{\mathfrak{m}}/K_{\mathfrak{m},1}\hookrightarrow\mathbb{I}/K^{\times}.
\]
至于满性, 应用弱逼近定理 (ANT, 7.20). 设 $S=S(\mathfrak{m})$, $\mathbf{a}=(a_v)\in\mathbb{I}$. 若取 $b\in K$ 使其在 $K_v^{\times}$ 中对所有 $v\in S$ 都非常接近 $a_v$, 则 $a_v/b$ 在 $K_v^{\times}$ 中对所有 $v\in S$ 接近 $1$; 事实上可取 $b$ 使 $a_v/b\in W_{\mathfrak{m}}(\mathfrak{p})$ 对所有 $v\in S$ 成立. 例如对 $S$ 中的实素 $v$, 只须取 $b$ 与 $a_v$ 在 $K_v$ 中同号. 此时 $\mathbf{a}/b\in\mathbb{I}_{\mathfrak{m}}$, 且它在 $\mathbb{I}/K^{\times}$ 中的像为 $\mathbf{a}$.
\end{proof}

\subsection*{理想的特征与 idèle 的特征}

设 $S\supset S_\infty$ 是 $K$ 的素的有限集, $G$ 是有限 Abel 群. 同态
\[
\psi\colon I^S\to G
\]
称为\textbf{容许一个模}, 若存在支撑含于 $S$ 的模 $\mathfrak{m}$ 使 $\psi(i(K_{\mathfrak{m},1}))=1$. 例如, 对每个 Abel 扩张 $L/K$, Artin 证明了 Artin 映射
\[
I^S\to\Gal(L/K)
\]
容许一个模.

\begin{proposition}\label{prop:v.4.7}
若 $\psi\colon I^S\to G$ 容许一个模, 则存在唯一的同态 $\phi\colon\mathbb{I}\to G$ 使得

(a) $\phi$ 连续 ($G$ 取离散拓扑);

(b) $\phi(K^{\times})=1$;

(c) $\phi(\mathbf{a})=\psi(\mathrm{id}(\mathbf{a}))$, 对所有 $\mathbf{a}\in\mathbb{I}^S\overset{\mathrm{def}}{=}\{\mathbf{a}\mid a_v=1\ \text{对所有}\ v\in S\}$.

此外, 每个满足 (b) 的连续同态 $\phi\colon\mathbb{I}\to G$ 都由某个 $\psi$ 产生.
\end{proposition}

\begin{proof}
因 $\psi$ 容许模 $\mathfrak{m}$, 它经过 $I^{\mathfrak{m}}/i(K_{\mathfrak{m},1})=C_{\mathfrak{m}}$ 分解. 于是有图:
\begin{equation}
\begin{tikzcd}[column sep=small]
I^{\mathfrak{m}}\ar[r] & C_{\mathfrak{m}}\ar[r, "\psi"] & G\\
\mathbb{I}_{\mathfrak{m}}/K_{\mathfrak{m},1}\ar[r]\ar[u, "\simeq"] & \mathbb{I}_{\mathfrak{m}}/K_{\mathfrak{m},1}W_{\mathfrak{m}}\ar[u, "\simeq"'] & \\
\mathbb{I}\ar[r]\ar[u, hookrightarrow] & \mathbb{I}/K^{\times}\ar[u, two heads] &
\end{tikzcd}\tag{43}
\end{equation}
其中的同构即命题 \ref{prop:v.4.6} 中的那些, 其余未命名的映射是商映射. 定义 $\phi$ 为复合 $\mathbb{I}\to G$. 它显然具有性质 (a) 与 (b), 且
\[
\phi(\mathbf{a})=\psi(\mathrm{id}(\mathbf{a}))\quad\text{对所有}\ \mathbf{a}\in\mathbb{I}_{\mathfrak{m}},
\]
从而更是满足性质 (c).

为证该映射由 (a), (b), (c) 唯一确定, 只须证 $\mathbb{I}^S K^{\times}$ 在 $\mathbb{I}$ 中稠密, 而这由弱逼近定理 (ANT, 7.20) 推出: 设 $\mathbf{a}\in\mathbb{I}$; 取 $b\in K^{\times}$ 使其对 $v\in S$ 非常接近 $a_v$, 并令 $\mathbf{a}'$ 为满足 $a_v'b=a_v$ (对所有 $v\notin S$) 的 $\mathbb{I}^S$ 中的元素. 则 $\mathbf{a}'b\in\mathbb{I}^S\cdot K$ 且在 $\mathbb{I}$ 中接近 $\mathbf{a}$.

反之, 设 $\phi\colon\mathbb{I}\to G$ 连续. 其核包含 $1$ 的一个开邻域, 故对某个 $S$ 与 $\varepsilon$ 有 $U(S,\varepsilon)\subset\mathrm{Ker}(\phi)$. 考虑无限素 $v$. $\phi$ 在 $K_v^{\times}$ 上的限制是连续映射 $\R^{\times}\to G$ 或 $\CC^{\times}\to G$. 显然 $K_v^{\times}$ 中含 $1$ 的连通分支, 即 $\R_{>0}$ 或 $\CC^{\times}$, 映到 $1$, 故在核中. 综合这些观察, 可见 $\phi$ 的核包含某个 $\mathfrak{m}$ 的 $W_{\mathfrak{m}}$.

现在可以再次使用证明开头的图. 我们有同态 $\phi\colon\mathbb{I}/K^{\times}\to G$, 可把它\guillemotleft 限制\guillemotright\ 为同态 $\mathbb{I}_{\mathfrak{m}}/K_{\mathfrak{m},1}\to G$. 该同态在 $W_{\mathfrak{m}}$ 上平凡, 故经过 $\mathbb{I}_{\mathfrak{m}}/K_{\mathfrak{m},1}W_{\mathfrak{m}}$ 分解. 此同态可转移到 $C_{\mathfrak{m}}$ 上, 再与 $\mathbb{I}\to C_{\mathfrak{m}}$ 复合. 这便是要找的 $\psi$.
\end{proof}

\begin{remark}\label{rem:v.4.8}
设 $G$ 为交换拓扑群. 称同态 $\psi\colon I^S\to G$ \textbf{可容许的} (admissible), 若对 $G$ 中 $1$ 的每个邻域 $N$, 存在模 $\mathfrak{m}$ 使 $\psi(i(K_{\mathfrak{m},1}))\subset N$. 则每个可容许同态 $\psi$ 定义满足命题 \ref{prop:v.4.7} 条件 (a), (b), (c) 的同态 $\phi\colon\mathbb{I}\to G$. 此外, 若 $G$ 完备且\guillemotleft 无小子群\guillemotright, 即存在 $1$ 的不含非平凡子群的邻域, 则每个满足 (b) 的连续同态 $\phi\colon\mathbb{I}\to G$ 都由某个可容许的 $\psi$ 产生. 证明与命题 \ref{prop:v.4.7} 相同 (见 \cite{Tate1967}, 命题 4.1).

圆群 $G=\{z\in\CC\mid|z|=1\}$ 完备且无小子群. 可容许的 $\psi\colon I^S\to G$ 及相应的 $\phi$ 称为 \textbf{Hecke 特征}\index{heckelezheng@Hecke 特征(Hecke character)}.
\end{remark}

\begin{remark}\label{rem:v.4.9}
给定 $\psi$ 后我们选取了一个 $\mathfrak{m}$, 然后说明如何构造 $\phi$. 实践中更常见的做法是: 在已知 $\phi$ 满足条件 (a), (b), (c) 的情况下直接把它认出来. 为此, 下列观察有用.

(a) 设 $\mathbf{a}=(a_v)$ 是满足 $a_v=1$ 对所有有限素、$a_v>0$ 对所有实素的 idèle; 则 $\phi(\mathbf{a})=1$. 为看出这一点, 注意 $\prod_{v|\infty}K_v^{\times}$ 作为 $\mathbb{I}$ 的子群上的诱导拓扑是其自然拓扑, 故 $\phi$ 在其上的限制在含 $1$ 的连通分支上平凡.

(b) 设 $\mathbf{a}=(a_v)$ 是满足 $a_v=1$ 对所有 $v\in S$、$a_v$ 为单位对所有 $v\notin S$ 的 idèle; 则 $\phi(\mathbf{a})=1$. 事实上这直接由条件 (c) 推出.

(c) 若 $\mathbf{a}$ \guillemotleft 接近 $1$\guillemotright, 则 $\phi(\mathbf{a})=1$. 事实上, 鉴于 $G$ 取离散拓扑, 这直接由条件 (a) 推出.

(d) 结合 (a), (b), (c) 可知: 若 $\mathbf{a}=(a_v)$ 满足
\[
\begin{cases}
a_v>0 & \text{当 } v\ \text{实};\\
a_v\ \text{接近}\ 1 & \text{当 } v\in S\ \text{有限};\\
a_v\ \text{为单位} & \text{当 } v\notin S,
\end{cases}
\]
则 $\phi(\mathbf{a})=1$. 事实上, (a) 与 (b) 表明可以用某些类型的 idèle 乘 $\mathbf{a}$ 而不改变 $\phi(\mathbf{a})$ 的值. 显然, 若 $a_v$ 对 $S$ 中的有限 $v$ 接近 $1$, 可以用这样的 idèle 乘之使其更接近 $1$.
\end{remark}

\begin{example}\label{ex:v.4.10}
设 $L=\Q[\zeta_p]$, $\psi$ 为 Artin 映射
\[
I^S\to(\Z/p\Z)^{\times}\to\Gal(L/\Q),\qquad S=\{p,\infty\}.
\]
回顾: 第一个映射把 (唯一) 由 $(r/s)$, $r,s>0$, $(p,r)=1=(p,s)$, 代表的理想送到 $[r][s]^{-1}$, 第二个映射把 $[m]$ 送到自同构 $\zeta\mapsto\zeta^m$. 总的说来, 对任意素数 $l\ne p$, 该映射把 $(l)$ 送到 $l$ 处的 Frobenius 自同构 $\zeta\mapsto\zeta^l$. 设 $\phi\colon\mathbb{I}\to\Gal(L/\Q)$ 为定理中对应于 $\psi$ 的同态. 我们要显式确定 $\phi$.

设 $\mathbf{a}=(a_\infty,a_2,\ldots,a_p,\ldots,a_l,\ldots)$ 是 $\Q$ 的 idèle. 若 $a_\infty=1=a_p$, 则 $\phi(\mathbf{a})=\psi(\mathrm{id}(\mathbf{a}))$. 于是 $\phi(\mathbf{a})=\zeta_p^m$, 其中 $m=\prod l^{\mathrm{ord}_l(a_l)}$.

考虑 $\mathbf{p}=(1,\ldots,1,p,1,\ldots)$ ($p$ 在第 $p$ 个位置). 则
\[
\mathbf{p}/p=(p^{-1},\ldots,p^{-1},1,p^{-1},\ldots).
\]
据上述注记的 (d), $\phi(\mathbf{p}/p)=1$, 故
\[
\phi(\mathbf{p})=\phi(\mathbf{p}/p)\phi(p)=1.
\]
此处 $p$ 既表示元素 $p\in\Q_p$, 也表示主 idèle $(p,p,\ldots)$.

考虑 $\mathbf{a}=(1,\ldots,1,u,1,\ldots)$, $u\in\Z_p^{\times}$, $u$ 在第 $p$ 个位置. 写
\[
u^{-1}=a_0+a_1p+\cdots+a_sp^s+\cdots,\qquad 0\le a_i<p,\quad a_i\in\Z,
\]
并令 $c=a_0+\cdots+a_sp^s\in\Z$. 注意 $c>0$. 则 $uc\in 1+p^{s+1}\Z_p$, 即对大的 $s$ 它\guillemotleft 接近 $1$\guillemotright. 写
\[
\mathbf{a}c=(c,c,\ldots,c,\overset{l|c}{1},c,\ldots,\overset{p}{uc},c,\ldots)\,(1,\ldots,1,\overset{l|c}{c},1,\ldots).
\]
第一个因子是 $\mathbf{a}c$, 只是把整除 $c$ 的素 $l$ 处的分量挪到了第二个因子. 对大的 $s$, 由上述注记的 (d) 有 $\phi(\text{第一个因子})=1$. 第二个因子属于 $\mathbb{I}^S$, 而我们对 $\phi|\mathbb{I}^S$ 的描述表明 $\phi(\text{第二个因子})$ 把 $\zeta$ 送到 $\zeta^c$. 结论:
\begin{equation}
\phi(\mathbf{a})(\zeta)=\zeta^c=\zeta^{u^{-1}}.\tag{44}
\end{equation}

考虑 $\mathbf{a}=(-1,1,\ldots,1)$. 则
\[
-\mathbf{a}=(1,-1,\ldots,-1,\overset{p}{1},-1,\ldots)\,(1,\ldots,1,\overset{p}{-1},1,\ldots),
\]
而
\[
\phi(\mathbf{a})=\phi(-\mathbf{a})=\phi(1,\ldots,1,-1,1,\ldots).
\]
据 (44), $\phi(\mathbf{a})(\zeta)=\zeta^{-1}$.

由于 $\phi$ 是同态, 这就完成了对它的显式描述.
\end{example}

\begin{remark}\label{rem:v.4.11}
由 (43) 得到典范同态 $\pi_{\mathfrak{m}}\colon\mathbb{I}\to C_{\mathfrak{m}}$. 这是唯一的连续同态 $\pi_{\mathfrak{m}}\colon\mathbb{I}\to C_{\mathfrak{m}}$ 使得

(a) $\pi_{\mathfrak{m}}(K^{\times})=1$;

(b) $\pi_{\mathfrak{m}}(\mathbf{a})=\mathrm{id}(\mathbf{a})$ 对所有 $\mathbf{a}\in\mathbb{I}^{S(\mathfrak{m})}$.

(在命题 \ref{prop:v.4.7} 中取 $\psi$ 为 $I^S\to C_{\mathfrak{m}}$.)

若 $\mathfrak{m}|\mathfrak{m}'$, 则 $\pi_{\mathfrak{m}'}$ 与典范同态 $C_{\mathfrak{m}'}\to C_{\mathfrak{m}}$ 的复合满足刻画 $\pi_{\mathfrak{m}}$ 的条件. 因此诸 $\pi_{\mathfrak{m}}$ 合成一个连续同态 $\pi\colon\mathbb{I}\to\varprojlim C_{\mathfrak{m}}$. 我们要确定该映射的核与像.

由于每个 $\pi_{\mathfrak{m}}\colon\mathbb{I}\to C_{\mathfrak{m}}$ 都是满射, 像是稠密的. 事实上 $\mathbb{I}^1\to C_{\mathfrak{m}}$ 是满射, 故 $\pi_{\mathfrak{m}}(\mathbb{I}^1)$ 稠密. 但 $\pi_{\mathfrak{m}}(\mathbb{I}^1)$ 是紧的, 因为 $\pi_{\mathfrak{m}}$ 经过紧群 $\mathbb{I}^1/K^{\times}$ 分解, 因而是完备的. 这表明 $\pi$ 是满射.

设 $\mathbb{I}_\infty^+$ 为满足 $a_v=1$ ($v$ 有限)、$a_v>0$ ($v$ 实) 的 idèle $\mathbf{a}$ 之集. 则 $\mathbb{I}_\infty^+$ 同构于 $(K\otimes_{\Q}\R)^{\times}=\prod_{v|\infty}K_v^{\times}$ 的单位分支. $\mathbb{I}\to\varprojlim C_{\mathfrak{m}}$ 的核包含 $\mathbb{I}_\infty^+\cdot K^{\times}$, 从而包含其闭包. 事实上它就等于闭包. [在将来的版本中, 将更详细地考察这些内容.]
\end{remark}

\subsection*{idèle 的范}

设 $L$ 是数域 $K$ 的有限扩张, $v$ 是 $K$ 的素. 回顾 (ANT, 8.2): 存在典范同构
\[
L\otimes_K K_v\to\prod_{w|v}L_w.
\]
由此 (同书 8.3) 对任意 $\alpha\in L$,
\[
\Nm_{L/K}\alpha=\prod_{w|v}\Nm_{L_w/K_v}\alpha\qquad(\text{等式在}\ K_v\ \text{内}).
\]
对 idèle $\mathbf{a}=(a_w)\in\mathbb{I}_L$, 定义 $\Nm_{L/K}(\mathbf{a})$ 为 $\mathbb{I}_K$ 中满足 $b_v=\prod_{w|v}\Nm_{L_w/K_v}a_w$ 的 idèle $\mathbf{b}$. 前面的注记表明下述图中左侧的正方形交换, 右侧的正方形显然也交换:
\[
\begin{tikzcd}
L^{\times}\ar[r]\ar[d, "\Nm_{L/K}"'] & \mathbb{I}_L\ar[r, "\mathrm{id}"]\ar[d, "\Nm_{L/K}"'] & I_L\ar[d, "\Nm_{L/K}"']\\
K^{\times}\ar[r] & \mathbb{I}_K\ar[r, "\mathrm{id}"'] & I_K
\end{tikzcd}
\]
于是得到交换图:
\[
\begin{tikzcd}
C_L\ar[r]\ar[d, "\Nm_{L/K}"'] & \mathbf{C}_L\ar[d, "\Nm_{L/K}"']\\
C_K\ar[r] & \mathbf{C}_K
\end{tikzcd}
\]
($C_K=$ 理想类群 $I/i(K^{\times})$; $\mathbf{C}_K=$ idèle 类群 $\mathbb{I}/K^{\times}$.)

\begin{proposition}\label{prop:v.4.12}
若 $L/K$ 是特征零的局部域的有限扩张, 则

(a) $\Nm_{L/K}(L^{\times})=\R_{>0}$ (情形 $K=\R$, $L=\CC$);

(b) $\Nm_{L/K}(L^{\times})\supset 1+\mathfrak{p}_K^m$ 对某个 $m$ (情形 $K$ 非阿基米德);

(c) $\Nm_{L/K}(L^{\times})\supset\mathcal{O}_K^{\times}$ (情形 $K$ 非阿基米德且 $L/K$ 非分歧).
\end{proposition}

\begin{proof}
(a) 显然. (b) 见 (I \ref{lem:1.3}), (c) 见 (III, \ref{prop:iii.1.2}).
\end{proof}

\begin{corollary}\label{cor:v.4.13}
设 $L/K$ 是数域的有限扩张. 则 $\Nm_{L/K}\mathbb{I}_L\supset W_{\mathfrak{m}}$ 对某个模 $\mathfrak{m}$.
\end{corollary}

\section{idèle 语言的主要定理}\label{sec:V.5}

用理想陈述的类域论主要定理非常显式, 且对许多目的而言是最有用的表述. 但它有一些缺点: 必须先固定一个模 $\mathfrak{m}$, 而该理论只描述导子整除 $\mathfrak{m}$ 的那些 Abel 扩张; 特别地, 它对 $K$ 的无限 Abel 扩张不作任何描述. 用 idèle 陈述的主要定理使人能同时考虑无限 Abel 扩张, 或等价地, 一并考虑所有有限 Abel 扩张. 它还使整体与局部 Artin 映射之间的关系变得透明.

设 $L$ 是 $K$ 的有限 Abel 扩张. 设 $v$ 是 $K$ 的素, $w$ 是 $L$ 中位于 $v$ 之上的素. 回顾 $w$ 的分解群 $D(w)$ 是子群
\[
D(w)=\{\sigma\in\Gal(L/K)\mid\sigma w=w\}.
\]
其元素唯一地延拓为 $L_w/K_v$ 的自同构, 且 $D(w)\simeq\Gal(L_w/K_v)$. 局部类域论 (I, \ref{thm:1.1}) 给出同态 (局部 Artin 映射)
\[
\phi_v\colon K_v^{\times}\to D(w)\subset G.
\]

\begin{lemma}\label{lem:v.5.1}
$G$ 的子群 $D(w)$ 与映射 $\phi_v$ 不依赖于位于 $v$ 之上的素 $w$ 的选取.
\end{lemma}

\begin{proof}
位于 $v$ 之上的其他素形如 $\sigma w$, $\sigma\in G$, 且 $\sigma\colon L\to L$ 由连续性延拓为保持 $K_v$ 的同态 $\sigma\colon L_w\to L_{\sigma w}$. 我们有
\[
D(\sigma w)=\sigma D(w)\sigma^{-1},
\]
而它等于 $D(w)$, 因为 $G$ 交换.

设 $\Omega$ 与 $\Omega'$ 分别是 $K_v$ 的包含 $L_w$ 与 $L_{\sigma w}$ 的极大 Abel 扩张. 由第 \ref{ch:III} 章, 我们有局部 Artin 映射 $\phi_v\colon K^{\times}\to\Gal(\Omega/K_v)$ 与 $\phi_v'\colon K^{\times}\to\Gal(\Omega'/K_v)$ 同构 $\tilde{\sigma}\colon\Omega\to\Omega'$ 的选取确定同构
\[
\rho\mapsto\tilde{\sigma}\circ\rho\circ\tilde{\sigma}^{-1}\colon\Gal(\Omega/K_v)\to\Gal(\Omega'/K_v),
\]
它不依赖于 $\tilde{\sigma}$. 此外它与 $\phi_v$ 的复合是 $\phi_v'$ (因为它满足刻画 $\phi_v'$ 的条件).
\end{proof}

\begin{proposition}\label{prop:v.5.2}
存在唯一的连续同态 $\phi_K\colon\mathbb{I}\to\Gal(K^{\mathrm{ab}}/K)$ 具有如下性质: 对任何含于 $K^{\mathrm{ab}}$ 的 $K$ 的有限扩张 $L$ 及 $L$ 中位于 $K$ 的素 $v$ 之上的素 $w$, 图
\[
\begin{tikzcd}
K_v^{\times}\ar[r, "\phi_v"]\ar[d] & \Gal(L_w/K_v)\ar[d, "\text{限制}\ |\!\ _L"]\\
\mathbb{I}_K\ar[r, "\mathbf{a}\mapsto\phi_K(\mathbf{a})|L"'] & \Gal(L/K)
\end{tikzcd}
\]
交换.
\end{proposition}

\begin{proof}
设 $\mathbf{a}\in\mathbb{I}$, $L\subset K^{\mathrm{ab}}$ 在 $K$ 上有限. 若 $a_v\in U_v$ 且 $L_w/K_v$ 非分歧, 则 $\phi_v(a_v)=1$ (见 III, \ref{sec:III.1}). 因此除有限多个 $v$ 外 $\phi_v(a_v)=1$, 故可定义
\[
\phi_{L/K}(\mathbf{a})=\prod_v\phi_v(a_v)
\]
(乘积在 $\Gal(L/K)$ 内). 显然 $\phi_{L/K}$ 是使上图交换的唯一同态.

若 $L'\supset L$, 则局部 Artin 映射的性质表明 $\phi_{L'/K}(\mathbf{a})|L=\phi_{L/K}(\mathbf{a})$. 因此存在唯一的同态 $\phi\colon\mathbb{I}\to\Gal(K^{\mathrm{ab}}/K)$ 使得 $\phi(\mathbf{a})|L=\phi_{L/K}(\mathbf{a})$ 对所有含于 $K^{\mathrm{ab}}$ 的 $K$ 的有限扩张 $L$ 成立.

再由局部 Artin 映射的性质, 对任何 $K\subset K'\subset L\subset K^{\mathrm{ab}}$ ($L$ 在 $K$ 上有限), 图
\[
\begin{tikzcd}
\mathbb{I}_{K'}^S\ar[r, "\phi_{L/K'}"]\ar[d, "\Nm"'] & \Gal(L/K')\ar[d]\\
\mathbb{I}_K^S\ar[r, "\phi_{L/K}"'] & \Gal(L/K)
\end{tikzcd}
\]
交换. 取 $K'=L$ 可知 $\Nm_{L/K}(\mathbb{I}_L^S)$ 含于 $\phi_{L/K}$ 的核. 特别地, $\phi_{L/K}$ 的核包含 $\mathbb{I}_K^S$ 的一个开子群 (推论 \ref{cor:v.4.13}), 这蕴含 $\phi_K$ 连续.
\end{proof}

\begin{theorem}[互反律]\label{thm:v.5.3}
同态 $\phi_K\colon\mathbb{I}_K\to\Gal(K^{\mathrm{ab}}/K)$ 具有如下性质:

(a) $\phi_K(K^{\times})=1$;

(b) 对 $K$ 的每个有限 Abel 扩张 $L$, $\phi_K$ 定义同构
\[
\phi_{L/K}\colon\mathbb{I}_K/(K^{\times}\cdot\Nm(\mathbb{I}_L))\to\Gal(L/K).
\]
\end{theorem}

我们在命题的证明中已看到 $\phi_{L/K}(\Nm(\mathbb{I}_L))=1$, 故 (假设定理的 (a)) $\phi_{L/K}$ 确实经过 $\mathbb{I}_K/K^{\times}\cdot\Nm(\mathbb{I}_L)$ 分解. (b) 也可表述为: $\phi$ 定义同构
\[
\phi_{L/K}\colon\mathbf{C}_K/\Nm(\mathbf{C}_L)\to\Gal(L/K).
\]

\begin{example}\label{ex:v.5.4}
定理的 (a) 说: 对每个 $b\in K^{\times}$, $\prod\phi_v(b)=1$. 把它应用于扩张 $K[a^{\frac{1}{n}}]/K$ (设 $K$ 含本原 $n$ 次单位根), 便得到 Hilbert 符号的乘积公式:
\[
\prod_v(a,b)_v=1.
\]
见 (III, \ref{rem:iii.4.8}).
\end{example}

\begin{theorem}[存在定理]\label{thm:v.5.5}
固定 $K$ 的一个代数闭包 $K^{\mathrm{al}}$; 对 $C_K$ 的每个有限指数开子群 $N$, 存在唯一的含于 $K^{\mathrm{al}}$ 的 $K$ 的 Abel 扩张 $L$ 使得 $\Nm_{L/K}C_L=N$.
\end{theorem}

$C_K$ 的子群称为\textbf{范群}, 若它形如 $\Nm(C_L)$, $L$ 为 $K$ 的某个有限 Abel 扩张. 存在定理表明: 范群恰好是 $C_K$ 的有限指数开子群. 若 $N$ 是这样的群, 则满足 $\Nm(C_L)=N$, 即 $N=\mathrm{Ker}(\phi_{L/K})$, 的 $K$ 的有限 Abel 扩张 $L$ 称为属于 $N$ 的 $K$ 的\textbf{类域}.

如上陈述的存在定理对所有整体域成立.

\begin{corollary}\label{cor:v.5.6}
映射 $L\mapsto\Nm(C_L)$ 是从 $K$ 的有限 Abel 扩张之集到 $C_K$ 的有限指数开子群之集的双射. 此外,
\begin{align*}
L_1\subset L_2&\iff\Nm(\mathbf{C}_{L_1})\supset\Nm(\mathbf{C}_{L_2});\\
\Nm(\mathbf{C}_{L_1\cdot L_2})&=\Nm(\mathbf{C}_{L_1})\cap\Nm(\mathbf{C}_{L_2});\\
\Nm(\mathbf{C}_{L_1\cap L_2})&=\Nm(\mathbf{C}_{L_1})\cdot\Nm(\mathbf{C}_{L_2}).
\end{align*}
\end{corollary}

\begin{remark}\label{rem:v.5.7}
\begin{enumerate}[(a)]
\item 在数域情形, 映射
\[
\phi_K\colon\mathbb{I}_K\to\Gal(K^{\mathrm{ab}}/K)
\]
是满射. 对 $K$ 的无限素 $v$, 记 $K_v^+$ 为 $K_v^{\times}$ 中含 $1$ 的连通分支; 于是依 $v$ 为复或实, $K_v^+$ 同构于 $\CC^{\times}$ 或 $\R_{>0}$. 显然 $\prod_{v|\infty}K_v^+\subset\mathrm{Ker}(\phi_K)$. 按定义 $K^{\times}\subset\mathrm{Ker}(\phi_K)$, 故 $K^{\times}\cdot(\prod_{v|\infty}K_v^+)\subset\mathrm{Ker}(\phi_K)$. 但 $\phi_K$ 是连续同态而 $\Gal(K^{\mathrm{ab}}/K)$ 是 Hausdorff 的, 故核是闭子群. 于是 $\mathrm{Ker}(\phi_K)$ 包含 $K^{\times}\cdot(\prod_{v|\infty}K_v^+)$ 的闭包. 这恰好就是核是一条定理. $K^{\times}\cdot(\prod_{v|\infty}K_v^+)$ 的闭包在 $C_K$ 中的像是 $C_K$ 中含 $1$ 的连通分支.

对 $K$ 的每个有限 Abel 扩张 $L$, Artin 映射定义同构 $C_K/\Nm(C_L)\to\Gal(L/K)$. 对 $L$ 取逆向极限, 右端得到与 $\Gal(K^{\mathrm{ab}}/K)$ 的同构, 于是问题化为计算体系 $C_K/\Nm(C_L)$ 的逆向极限. 存在定理表明群 $\Nm(C_L)$ 恰为有限指数开子群, 故这是一个拓扑学问题.

\item 在函数域情形, Artin 映射 $\phi_K\colon\mathbb{I}_K/K^{\times}\to\Gal(K^{\mathrm{ab}}/K)$ 是单射, 但不满 (其像稠密).
\end{enumerate}
\end{remark}

\begin{remark}\label{rem:v.5.8}
设整体 Artin 映射 $\phi\colon\mathbb{I}\to\Gal(K^{\mathrm{ab}}/K)$ 的核包含 $K^{\times}$. 则对 $K$ 的每个有限 Abel 扩张 $L/K$, $\phi_{L/K}\colon\mathbb{I}\to\Gal(L/K)$ 由一个容许模的同态 $\psi\colon I^S\to\Gal(L/K)$ 产生 (如命题 \ref{prop:v.4.7}). 此外, 由于 $\phi$ 是局部 Artin 映射之积, $\psi$ 必为理想语言的整体 Artin 映射 (因而它容许一个模). 由 idèle 版本的定理\ref{thm:v.5.3} 与 \ref{thm:v.5.5} 推出理想版本的定理\ref{thm:v.3.5} 与 \ref{thm:v.3.6} 是一个直接的练习. 我们将在第 \ref{ch:VII} 章证明定理\ref{thm:v.5.3} 与 \ref{thm:v.5.5}.
\end{remark}

\phantomsection\label{pg:orig180}%
\begin{NOTES}
在其 1951 年的报告 Sur la théorie du corps de classes 中, Weil 写道:\footnote{寻找 $C_k$ 的一个解释 ($k$ 为数域时), 使其在某种意义上类似于 $k$ 为函数域时由 Galois 群给出的解释, 在我看来是当前数论的基本问题之一; 也许这样的解释蕴含着 Riemann 假设的钥匙\ldots\ (原书脚注给出的英译, 此处译出.)}
\begin{quote}
\itshape La recherche d'une interprétation de $C_k$ si $k$ est un corps de nombres, analogue en quelque manière \`a l'interprétation par un groupe de Galois quand $k$ est un corps de fonctions, me semble constituer l'un des problèmes fondamentaux de la théorie des nombres \`a l'heure actuelle; il se peut qu'une telle interprétation renferme la clef de l'hypothèse de Riemann\ldots
\end{quote}
\end{NOTES}
关于此段的讨论见 \mo{41296}.

\subsection*{例}

\begin{lemma}\label{lem:v.5.9}
映射
\[
(r,t,(u_p))\mapsto(rt,ru_2,ru_3,ru_5,\ldots)\colon\Q^{\times}\times\R_{>0}\times\prod\Z_p^{\times}\to\mathbb{I}_{\Q}
\]
是拓扑群的同构 ($\Q^{\times}$ 取离散拓扑).
\end{lemma}

\begin{proof}
任何 idèle $\mathbf{a}=(a_\infty,a_2,\ldots,a_p,\ldots)$ 可写成
\[
\mathbf{a}=a(t,u_2,u_3,u_5,\ldots),\qquad a\in\Q^{\times},\ t\in\R_{>0},\ u_p\in\Z_p^{\times},
\]
\,---\,取 $a=(\mathrm{sign}(a_\infty))\prod p^{\mathrm{ord}_p(a_p)}$, $t=a_\infty/a$, $u_p=a_p/a$. 而且表达式唯一, 因为对所有的 $p$ 都是 $p$-adic 单位的正有理数只有 $1$.

形如
\[
\{1\}\times U\times\prod_{p\ \mathrm{finite}}U_p
\]
的集合 ($U$, $U_p$ 分别为 $1$ 在 $\R^{\times}$, $\Q_p^{\times}$ 中的开邻域, 且除有限多个 $p$ 外 $U_p=\Z_p^{\times}$) 在左边构成 $1$ 的基本邻域组, 在右边亦然.
\end{proof}

于是有拓扑群的典范同构
\[
C_{\Q}\to\R_{>0}\times\prod_{p\ \mathrm{finite}}\Z_p^{\times}=\R_{>0}\times\hat{\Z}^{\times}.
\]
令
\[
\Q^{\mathrm{cyc}}=\bigcup\Q[\zeta_n].
\]
此情形下, 整体互反映射是
\[
\phi\colon\mathbb{I}_{\Q}\to\hat{\Z}^{\times}\to\Gal(\Q^{\mathrm{cyc}}/\Q)
\]
的\guillemotleft 逆\guillemotright, 其中 $\mathbb{I}_{\Q}\to\hat{\Z}^{\times}$ 是上述投影映射, 而 $\hat{\Z}^{\times}\to\Gal(\Q^{\mathrm{cyc}}/\Q)$ 是典范同构 (见 I A.5c).

\begin{summary}\label{note:v.5.10}
\phantomsection\hypertarget{item:v.5.10}{}
设 $K$ 是代数数域. 存在连续满同态 (互反映射或 Artin 映射)
\[
\phi_K\colon\mathbb{I}_K\to\Gal(K^{\mathrm{ab}}/K)
\]
使得对每个含于 $K^{\mathrm{ab}}$ 的 $K$ 的有限扩张 $L$, $\phi_K$ 给出交换图
\[
\begin{tikzcd}
\mathbb{I}_K/K^{\times}\ar[r, "\phi_K"]\ar[d] & \Gal(K^{\mathrm{ab}}/K)\ar[d, "\tau\mapsto\tau|L"]\\
\mathbb{I}_K/(K^{\times}\cdot\Nm(\mathbb{I}_L))\ar[r, "\phi_{L/K}"', "\simeq"] & \Gal(L/K)
\end{tikzcd}
\]
它由下列两条性质确定:

(a) 对每个满足如下条件的 $u=(u_v)\in\mathbb{I}_K$, $\phi_{L/K}(u)=1$:

\quad i) 若 $v$ 在 $L$ 中不分歧, 则 $u_v$ 是单位;

\quad ii) 若 $v$ 在 $L$ 中分歧, 则 $u_v$ 充分接近 $1$ (只依赖于 $L/K$);

\quad iii) 若 $v$ 是实素但在 $L$ 中变复, 则 $u_v>0$.

(b) 对 $K$ 中在 $L$ 不分歧的每个素 $v$, idèle
\[
\alpha=(1,\ldots,1,\overset{v}{\pi},1,\ldots),\qquad \pi\ \text{为}\ \mathcal{O}_v\ \text{的素元},
\]
映到 $\Gal(L/K)$ 中的 Frobenius 元 $(\mathfrak{p}_v,L/K)$.

为看出满足这些条件的映射至多一个, 设 $\alpha\in\mathbb{I}_K$, 用弱逼近定理取 $a\in K^{\times}$, 使其对所有在 $L$ 分歧或变复的素 $v$ 接近 $\alpha_v$. 则 $\alpha=au\beta$, 其中 $u$ 是 (a) 中的 idèle, $\beta$ 是形如 (b) 的 idèle 的有限乘积. 此时 $\phi_{L/K}(\alpha)=\phi_{L/K}(\beta)$, 后者可用 (b) 计算.

对 $K=\Q$, Artin 映射经过 $\{\pm\}\times\mathbb{I}_f/\Q^{\times}$ 分解, 而该商群的每个元素都被 $\hat{\Z}^{\times}\subset\mathbb{I}_f$ 中的元素唯一代表. 此情形下得到图
\[
\begin{tikzcd}
\hat{\Z}^{\times}\ar[r, "\phi_{\Q}", "\simeq"'] & \Gal(\Q^{\mathrm{ab}}/\Q) & =\bigcup_m\Q[\zeta_m]\\
(\Z/m\Z)^{\times}\ar[r, "\simeq"']\ar[u] & \Gal(\Q[\zeta_m]/\Q)\ar[u, hookrightarrow]
\end{tikzcd}
\]
(下方的同构把 $[n]$ 送到自同构 $\zeta_m\mapsto\zeta_m^n$.)
它交换, 且有逆. 这可以通过把 idèle $\alpha$ 写成上述形状 $au\beta$ 来验证, 但看一个例子更有启发. 设 $p$ 是不整除 $m$ 的素数, 并令
\[
\alpha=p\cdot(1,\ldots,1,p_p^{-1},1,\ldots)\in\Z\cdot\mathbb{I}_f=\mathbb{I}_f.
\]
则
\[
\alpha=(p,\ldots,p,\overset{p}{1},p,\ldots)\in\hat{\Z}^{\times}
\]
在 $(\Z/m\Z)^{\times}$ 中的像为 $[p]$, 而 $[p]$ 作为 $(p,\Q[\zeta_m]/\Q)$ 作用在 $\Q[\zeta_m]$ 上. 另一方面, $\phi_{\Q}(\alpha)=\phi_{\Q}((1,\ldots,1,p_p^{-1},1,\ldots))$, 它作为 $(p,\Q[\zeta_m]/\Q)^{-1}$ 作用.

\end{summary}
\begin{exercise}\label{exc:v.5.11}
证明: 即使 $K$ 是数域, $K$ 的 idèle 类群也有有限指数但不开的子群.
\end{exercise}

